\documentclass[11pt]{article}

\usepackage{amssymb,amsmath,amsfonts,amsthm}
\usepackage{latexsym}
\usepackage{graphics}
\usepackage{indentfirst}
\usepackage{lmodern}
\usepackage{hyperref}
\usepackage{mathrsfs}
\usepackage{array}
\usepackage{mathtools}
\usepackage{graphicx}
\usepackage{setspace}
\usepackage{enumerate}
\usepackage{tikz}
\usepackage{url,color}
\usepackage{xcolor}
\usepackage{chngcntr}
\usepackage{multicol}
\usepackage{comment}
\usepackage{float}
\usepackage{algpseudocode}
\usepackage{subcaption}
\usepackage[ruled]{algorithm2e}
\usepackage{booktabs}
\usepackage{caption}
\usepackage[letterpaper,margin=2cm,top=2.5cm,bottom=2.5cm]{geometry}

\AtBeginEnvironment{algorithm}{\setstretch{1.2}}

\hypersetup{colorlinks=true,citecolor=blue,linkcolor=blue,urlcolor=blue}

%%%%%%%%%%%%%%%%%%%%%%%%%%%%%%%%%%%%%%%%%%%%%%%%%%%%%%%%%%
%% Theorem environments
%%%%%%%%%%%%%%%%%%%%%%%%%%%%%%%%%%%%%%%%%%%%%%%%%%%%%%%%%%
\newtheorem{theorem}{Theorem}[section]
\newtheorem{lemma}{Lemma}[section]
\newtheorem{proposition}{Proposition}[section]
\newtheorem{corollary}{Corollary}[section]
\newtheorem{claim}{Claim}[section]
\newtheorem{conjecture}{Conjecture}[section]
\newtheorem{hypothesis}{Hypothesis}[section]
\newtheorem{assumption}{Assumption}[section]
\newtheorem{remark}{Remark}[section]
\newtheorem{example}{Example}[section]
\newtheorem{problem}{Problem}[section]
\newtheorem{definition}{Definition}[section]
\newtheorem{question}{Question}[section]
\newtheorem{condition}{Condition}[section]

%%%%%%%%%%%%%%%%%%%%%%%%%%%%%%%%%%%%%%%%%%%%%%%%%%%%%%%%%%
%% Math shorthands
%%%%%%%%%%%%%%%%%%%%%%%%%%%%%%%%%%%%%%%%%%%%%%%%%%%%%%%%%%
\newcommand{\N}{\mathbb{N}}
\newcommand{\Z}{\mathbb{Z}}
\newcommand{\E}{\mathbb{E}}
\newcommand{\Q}{\mathbb{Q}}
\newcommand{\R}{\mathbb{R}}
\newcommand{\C}{\mathbb{C}}
\newcommand{\norm}[1]{\left\lVert#1\right\rVert}

\allowdisplaybreaks

%%%%%%%%%%%%%%%%%%%%%%%%%%%%%%%%%%%%%%%%%%%%%%%%%%%%%%%%%%
%% Title, authors, abstract
%%%%%%%%%%%%%%%%%%%%%%%%%%%%%%%%%%%%%%%%%%%%%%%%%%%%%%%%%%

\title{Diffusion Proposals for Redistricting Markov Chains:\\
       Move-Size Theory, Exact Correction, and Ground-Truth Calibration}

\author{Alaittin K\i rt\i \c{s}o\u{g}lu \\
        \small Department of Applied Mathematics, Illinois Institute of Technology \\
        \small Chicago, IL 60616 \\
        \small \texttt{kirtisoglu@hawk.iit.edu}}

\date{\today}

\begin{document}

\maketitle

\begin{abstract}
Recombination (ReCom) chains are the standard Markov chain Monte
Carlo tool for sampling redistricting plans, and they are routinely
run under a Metropolis filter that nominally targets a
compactness-tempered distribution. We introduce a proposal built from
graph diffusion. One convex energy on the merged pair of districts, a
kernel-weighted Laplacian smoothness term plus a fidelity term that
anchors each move to the current plan, is minimized under global
population balance and rounded at random. The diffusion view yields a
small theory of the proposal. Its optimum is a single linear solve
whenever balance is slack, the expected number of units a proposal
moves is an explicit high-pass filter norm with a provable ceiling,
and the chain's mobility is governed by feasibility rather than by the
Metropolis ratio, which gives a tuning rule no existing redistricting
proposal has. Correcting the randomized rounding with an exact
Metropolis--Hastings step gives, to our knowledge, the first
optimization-based redistricting proposal with a detailed-balance
proof, verified to machine precision and by a batch-means test on an
enumerated fiber. On a fully enumerated $5 \times 5$ test bed we show
at the level of distributions rather than means that plain and
tempered ReCom do not sample their nominal target and sit close to
the spanning-tree-weighted measure, while the corrected chain samples
what it states. The same energy solved as an integer program gives an
anchored refinement tool, which we benchmark against annealing and
report as an optimizer rather than a sampler.
\end{abstract}

%════════════════════════════════════════════════════════════════════════════════
\section{Introduction}\label{sec:intro}
%════════════════════════════════════════════════════════════════════════════════

\subsection{The redistricting problem and ReCom's role}\label{sec:intro-problem}

Political redistricting seeks a partition of a geographic region into
electoral districts satisfying legal, demographic, and geometric
requirements, namely contiguity, population balance, compactness, and
preservation of communities of interest. Computational approaches model
the problem on the dual adjacency graph $G = (V,E)$ of geographic units
(precincts, census blocks, counties). A valid $k$-district plan is a
partition
\[
  P = \{V_1, \ldots, V_k\}, \qquad \bigsqcup_{i=1}^k V_i = V,
\]
where each $G[V_i]$ is connected and the populations $p(V_i)$ are
approximately equal. Let $\mathcal{P}_k(G)$ denote the space of valid
$k$-district plans. The cardinality of $\mathcal{P}_k(G)$ is enormous.
Even four congressional districts over the 99 counties of Iowa admit
roughly $10^{24}$ plans~\cite{fifield2020essential}.

Markov chain Monte Carlo (MCMC) has become the standard tool for
sampling $\mathcal{P}_k(G)$. Early flip chains modify a single
precinct's assignment per step~\cite{fifield2020automated} and mix
slowly. The Recombination (ReCom) algorithm of DeFord, Duchin, and
Solomon~\cite{deford2021recombination} performs multi-precinct updates
by merging two adjacent districts, sampling a uniform random spanning
tree of the merged subgraph, and cutting an edge to recover a balanced
bipartition. ReCom dramatically improved the mixing efficiency of
redistricting MCMC and is now the workhorse of the field, with
reversible~\cite{cannon2022spanning} and multi-scale~\cite{autry2020multiscale}
extensions.

\textbf{ReCom's blind spot.} The uniform spanning-tree split treats
every edge of the merged subgraph identically. The proposed
bipartition is determined by the topology of the merged subgraph
alone. Real geography is not topology-only. Shared boundary
length, population density, demographic similarity, and
administrative boundaries all carry signal that could shape the
proposed split. Practitioners have felt this. County-weighted
trees~\cite{clelland2022colorado} and GerryChain's region-aware
surcharges~\cite{gerrychain2024} inject indicator-based edge weights
into the tree draw, and SpecReCom~\cite{davies2025spectral} replaces
the tree with a Fiedler-vector cut of a randomly weighted Laplacian.
But these mechanisms are restricted (region indicators, hard-coded
county kernels, or no kernel at all), none includes a
locality-preserving fidelity term, and, as their authors
acknowledge, each shifts the chain's stationary distribution in
an uncharacterized, uncorrected way.

\subsection{Contribution}\label{sec:intro-contribution}

We introduce a kernel-aware local-refinement framework for ReCom-style
MCMC. The core object is a single Laplacian energy on the merged-pair
subgraph $H = G[V_i \cup V_j]$,
\begin{equation}\label{eq:core-energy}
  E(x) \;=\; \alpha\, x^\top L\, x \;+\; \beta\, \|x - x_0\|^2,
  \qquad
  (1 - \varepsilon)\bar p \;\le\; p^\top x \;\le\; (1+\varepsilon)\bar p,
\end{equation}
where $L$ is the kernel-weighted graph Laplacian of $H$, $x_0$ encodes
the current bipartition, and the balance constraint is anchored to the
global per-district ideal (Equation~\eqref{eq:balance-window}). The
Laplacian term $x^\top L x$ counts (or approximates) cut edges between
the two sides; the fidelity term $\|x - x_0\|^2$ penalizes deviation
from the current assignment, enforcing locality of each move. The
paper makes four contributions.

\begin{enumerate}
  \item \textbf{A diffusion proposal with a theory of its own
    behavior} (\S\ref{sec:qp-model}--\S\ref{sec:diffusion-theory}).
    The interior optimum of the energy is one implicit Euler step of
    the graph heat equation with timestep $\tau = \alpha/\beta$. We
    show that this step is a single linear solve whose output lies in
    the unit cube automatically, so the box constraints never bind
    and the proposal costs a linear solve whenever population balance
    is slack (Lemma~\ref{lem:row-stochastic}). Under single-shot
    Bernoulli rounding, the expected number of units a proposal moves
    is exactly the one-norm of a high-pass filter applied to the
    current split (Proposition~\ref{prop:move-size}). That quantity
    grows like twice the timestep times the pair's boundary for small
    timesteps and is bounded above by half the merged pair for any
    timestep (Corollary~\ref{cor:small-tau}), so mobility cannot be
    bought by tuning. We verify each statement against the solver
    and measure the missing factor, acceptance, directly. Given a
    feasible proposal the Metropolis step accepts it at least 98
    percent of the time; the chain's stickiness is entirely a
    feasibility effect with an interior optimum in $\tau$
    (\S\ref{sec:mobility}). To our knowledge no other redistricting
    proposal comes with a rule for sizing its moves.

  \item \textbf{An exactly reversible correction}
    (\S\ref{sec:qp-mh}). Correcting the rounding with an exact
    Metropolis--Hastings step yields Chain E, to our knowledge the
    first optimization-based redistricting proposal with a
    detailed-balance proof (Proposition~\ref{prop:qpmh-db}), stated
    on labeled plans, where the proposal map is a bijection. The
    proof identifies global anchoring of the balance window as the
    property that keeps forward and reverse solves symmetric. We
    verify the implemented kernel to machine precision on an
    enumerated fiber (Remark~\ref{rem:qpmh-scope}) and, on the same
    fiber, show by a batch-means test that a long run's occupancy
    distribution is not distinguishable from the exact target
    (\S\ref{sec:exp-2x3}). The ceiling of contribution 1 is what
    makes a global move necessary for mobility, and
    \S\ref{sec:disc-future} states the reversible hybrid this
    implies.

  \item \textbf{A ground-truth calibration of what the ReCom family
    samples, at the level of distributions}
    (\S\ref{sec:exp-5x5}--\S\ref{sec:exp-longrun}). Clelland et
    al.~\cite{clelland2021compactness} observed that ReCom's induced
    law decays roughly exponentially in cut length, and Procaccia and
    Tucker-Foltz~\cite{procaccia2022compact} proved that compact
    plans have many spanning trees. On a fully enumerated
    $5 \times 5$ fiber we make this exact. Using the Matrix--Tree
    theorem we compute the spanning-tree-weighted measure state by
    state and compare full empirical distributions against it and
    against the nominal tempered target, with sampling-noise floors
    at the effective sample size. The tempered chain's cut
    distribution is 220 floor standard deviations from its nominal
    target and close to, though distinguishable from, the
    tree-weighted tempered measure, while the exact chain is
    indistinguishable from the nominal target and far from the tree
    measure. The two chains separate exactly as theory predicts. All
    comparisons use effective-sample-size corrected inference and
    per-step balance audits, and we document two balance-constraint
    pitfalls, one affecting bare GerryChain \texttt{recom} proposals
    used without a validator, that materially change measured
    results.

  \item \textbf{The same energy as an anchored integer refinement}
    (\S\ref{sec:iqp}, \S\ref{sec:exp-compactness}--\S\ref{sec:exp-leastchange}).
    Solved as a convex MIQP inside a tempered chain, the energy gives
    Chain C, a \emph{matheuristic} ensemble generator in the lineage
    of optimization-based
    recombination~\cite{brous2025recom,cannon2023voting}, not a
    sampler (\S\ref{sec:chain-c-honest}). It reduces mean cut edges
    by $20.3\%$ versus Tempered ReCom on Iowa across ten seeds, and
    a simulated-annealing baseline we report matches its best plan
    at a fraction of the cost, so its value is anchored, kernel-aware
    refinement near a reference plan rather than compactness search.
    In a least-change experiment from Iowa's enacted map it ties the
    best annealer at a ten-county budget and is the only
    deterministic method there. The kernel showcase
    (\S\ref{sec:exp-multikernel}) and a structural diagnostic of the
    partition fiber (Appendix~\ref{sec:fiber-corridors}) support the
    main contributions.
\end{enumerate}

\paragraph{Positioning.} The construction of Chain E, a deterministic
optimization that sets the parameters of a random proposal which is
then corrected exactly, has ancestors outside redistricting.
Randomize-then-optimize~\cite{bardsley2014randomize} samples
posteriors by perturbing and solving, and Zanella's informed
proposals~\cite{zanella2020informed} give the general theory of
locally balanced, target-aware proposals on discrete spaces. Chain E
is an instance of that program on the space of plans, and our
priority claim is scoped to redistricting accordingly. Within
redistricting, the known-target samplers, Metropolized forest
recombination~\cite{autry2023metropolized}, RevReCom, Cycle Walk,
Marked Edge Walk, and
BUD~\cite{cannon2022spanning,deford2025cycle,mcwhorter2025marked,
akitaya2026bud}, achieve computable transition probabilities by
lifting to spanning-forest state spaces, while Chain E achieves the
same property at the partition level; the two routes are
complementary, and their global moves are the mobile component that
our exact-but-local kernel provably lacks
(Corollary~\ref{cor:small-tau}). Bonsai
sampling~\cite{clelland2026bonsai} draws independent plans from a
specified law without a chain and is the natural non-MCMC comparison
for any known-target method. \S\ref{sec:related} positions the
proposal against the weighted-tree, spectral, and matheuristic lines.

\paragraph{Roadmap.}\label{sec:intro-roadmap}
\S\ref{sec:related} reviews related work and \S\ref{sec:background}
sets up notation. \S\ref{sec:model} develops the proposal, its
theory, and the exact chain; \S\ref{sec:experiments} reports the
calibration, the convergence validation, and the Iowa, North
Carolina, kernel, and least-change experiments. Appendices hold
hyperparameters, the partition-fiber diagnostic, ablations, and a
reproducibility checklist.

%════════════════════════════════════════════════════════════════════════════════
\section{Related work}\label{sec:related}
%════════════════════════════════════════════════════════════════════════════════

\paragraph{MCMC for redistricting.}\label{sec:related-mcmc}

Early MCMC for redistricting framed the problem as a graph cut and used
single-unit flip chains~\cite{fifield2020automated} or
Swendsen--Wang-style large-swap moves. Flip chains mix poorly and
produce highly irregular districts~\cite{deford2021recombination,
najt2024empirical}. ReCom~\cite{deford2021recombination} introduced
the spanning-tree merge-split move and is now standard. Extensions
include Multi-Scale Merge--Split~\cite{autry2020multiscale} preserving
nested communities (counties, precincts), Forest ReCom~\cite{autry2023metropolized},
Sequential Monte Carlo~\cite{mccartan2023sequential} and its
generalized form~\cite{osullivan2026generalized}, whose repetition
effects are analyzed in~\cite{cannon2024repetition}, multiscale
parallel tempering~\cite{chuang2025multiscale}, and reversible
recombination~\cite{cannon2022spanning} targeting the canonical
spanning-tree distribution. The compactness-tempered target that
practitioners run under a Metropolis filter, and that we audit here,
originates in the North Carolina analyses of Herschlag et
al.~\cite{herschlag2020quantifying}. The 50-State Redistricting
Atlas~\cite{mccartan2022simulated} is the field's standard ensemble
benchmark. A consistent theme of the 2022--2026 literature is that a
sampler is only as defensible as its target.
RevReCom~\cite{cannon2022spanning}, Metropolized forest
recombination~\cite{autry2023metropolized}, the cycle
walk~\cite{deford2025cycle}, and the marked edge
walk~\cite{mcwhorter2025marked} all pay a structural price (lifted
state spaces, spanning-forest counting) precisely to keep transition
probabilities computable. On the theory side, Cannon, Pegden, and
Tucker-Foltz~\cite{cannon2024sampling} give a polynomial-time sampler
for the spanning-tree measure on balanced partitions of grids with a
constant number of districts, Gold~\cite{gold2026adversarial}
exhibits planar-grid configurations from which recombination needs
exponentially many expected steps to re-split, and mixing-time bounds
for the local chains used in practice remain open. Akitaya et
al.~\cite{akitaya2022reconfiguration} prove that deciding
reachability between balanced partitions under recombination is
PSPACE-complete when the district count and step size grow with the
graph.

\paragraph{Weighted and geography-aware proposals.}\label{sec:related-weighted}

Edge weighting inside the ReCom proposal is not new, and we position
our work against it explicitly. Clelland et
al.~\cite{clelland2022colorado} replace the proposal's edge-weighting
function with \emph{county-weighted trees} (intra-county edge
weights drawn from the uniform distribution $\mathcal{U}[0, w]$ for a tunable scalar $w$) to reduce county splits in Colorado ensembles. GerryChain's
region-aware ReCom~\cite{gerrychain2024} implements the same idea.
User-specified surcharges are added to the randomized Kruskal weights
for edges crossing designated region boundaries. Tapp~\cite{tapp2025compactness}
modifies the ReCom tree-and-cut step to improve shape-based
compactness at the proposal level. In all three, the weighting is an
indicator- or attribute-indexed heuristic, and (as their authors
acknowledge) the induced stationary distribution is uncharacterized
and uncorrected. What the unweighted tree proposal itself induces is
better understood. Clelland et al.~\cite{clelland2021compactness}
measured ReCom ensembles on grids and precinct graphs and found the
sampling probability of a plan to decay roughly exponentially in its
cut length, and Procaccia and
Tucker-Foltz~\cite{procaccia2022compact} proved the underlying fact
that compact plans have many spanning trees. Tapp, Proebsting, and
Ramsay~\cite{tapp2025parameter} study how tolerance and
county-preservation settings move ReCom ensembles and propose
effective-sample-size diagnostics for them. Our calibration in
\S\ref{sec:exp-5x5} makes the first of these observations exact on an
enumerated fiber, at the level of full distributions. Conversely, the
samplers with known targets cited above handle only unweighted
spanning-tree/forest measures or other tree-count-based families. The gap this paper addresses sits at the
intersection, namely a proposal family that accepts an \emph{arbitrary}
nonnegative edge kernel as a first-class input, together with (for the
Metropolis-corrected variant of \S\ref{sec:qp-mh}) an explicit
accept/reject step against a stated target.

\paragraph{Optimization-based districting and matheuristics.}\label{sec:related-opt}

A parallel operations research (OR) literature treats districting as
exact optimization. Examples include the Hess model lineage,
column-generation and set-partitioning
formulations~\cite{gurnee2021fairmandering}, multilevel
multiobjective methods~\cite{swamy2023multiobjective},
contiguity-constrained MIPs~\cite{validi2022imposing}, and cut-edge
minimization solved to optimality at county/tract
scale~\cite{validi2022political}. Closest to our least-change
experiment (\S\ref{sec:exp-leastchange}), Dobbs et
al.~\cite{dobbs2024facilitating} solve for minimal-change and
midpoint plans under transfer distance to reference plans exactly by
mixed-integer programming, which is the natural exact benchmark for
any anchored refinement. Between the sampling and optimization
families~\cite{becker2022redistricting}, hybrids are emerging. Short bursts~\cite{cannon2023voting} use ReCom
as a stochastic optimizer; Brous and Shmoys~\cite{brous2025recom}
replace the spanning-tree re-split with a Gurobi integer program
inside an accept-if-better local search for majority-minority
districts. Our integer chain (\S\ref{sec:iqp}) is closest to this
matheuristic line (it embeds convex MIQP local refinements inside a
recombination chain) and we present it as an
optimization-flavored ensemble generator, not a sampler with a
stationarity guarantee (\S\ref{sec:chain-c-honest}). The proximal
term $\beta\|x - x_0\|^2$ is the quadratic analogue of local
branching~\cite{fischetti2003local}, imposing a soft trust region
around the incumbent.

\paragraph{Spectral partitioning and diffuse-interface models.}\label{sec:related-spectral}

Spectral graph partitioning has a long history in image segmentation
and clustering~\cite{shi2000normalized}. In redistricting, the recent
SpecReCom of Davies, Job, Kampbell, Kim, and
Seo~\cite{davies2025spectral} (arXiv preprint, June 2025) replaces the
random spanning tree with a Fiedler-vector sign cut on a randomly
edge-weighted Laplacian of the merged pair, reporting roughly half the
cut edges of RevReCom on a $56 \times 56$ grid and a Colorado planar
graph. Plain SpecReCom sacrifices population balance, which their
BalSpecReCom variant partially restores. Diffuse-interface models
build on the Ginzburg--Landau (GL) functional, combining a smoothness
term with a fidelity term; the Chan--Vese model~\cite{chan2001active}
and Merriman--Bence--Osher (MBO) threshold dynamics implement this
iteratively. Bertozzi and collaborators~\cite{bertozzi2012diffuse,
merkurjev2013mbo} port the machinery to weighted graphs, where the
smoothness term becomes the Laplacian quadratic form $x^\top L x$ and
graph Tikhonov regularization~\cite{pilavci2021graph,shuman2013emerging}
recovers heat-kernel smoothing. Jacobs and Walch~\cite{jacobs2018pde}
brought threshold dynamics to redistricting before SpecReCom, using
volume-constrained MBO via auction dynamics~\cite{jacobs2018auction}
to generate compact equal-population maps. Theirs is a continuum PDE
approach, not a graph-kernel one, and it is not embedded in MCMC. To our
knowledge the present work is the first to use the \emph{graph} GL
lineage (kernel-weighted Laplacian smoothing plus fidelity) as a
proposal inside redistricting MCMC.

\paragraph{Cut minimization and discrete compactness.}\label{sec:related-cut}

The minimum graph cut~\cite{stoer1997simple,karger2000minimum} and its
balanced variants such as normalized cut~\cite{shi2000normalized} are
classical objectives. NP-hardness of balanced minimum
cut~\cite{wagner1993between} motivates continuous relaxations, of
which the Fiedler-vector cut and the Total Variation Isoperimetric
Profile~\cite{deford2019total} are two well-known examples. The cut
edge count $|\partial P|$ has been proposed as a discrete compactness
measure for redistricting~\cite{duchin2018discrete} and shown to
correlate well with classical geometric scores (Polsby--Popper, Reock)
while being more stable under unit refinement.

%════════════════════════════════════════════════════════════════════════════════
\section{Background: plans and ReCom}\label{sec:background}
%════════════════════════════════════════════════════════════════════════════════

\subsection{Dual graph and \texorpdfstring{$k$}{k}-district plans}\label{sec:background-graph}

Let $G = (V, E)$ be the dual graph of a state, with vertices
representing geographic units (precincts, census blocks, counties) and
edges encoding shared boundaries. Each vertex $v$ carries a population
$p_v \ge 0$ and may carry additional attributes (demographic vectors,
shapefile geometry). A \emph{$k$-district plan} is a partition
$P = \{V_1, \ldots, V_k\}$ of $V$ such that each induced subgraph
$G[V_i]$ is connected (\emph{contiguity}) and the populations satisfy
\begin{equation}\label{eq:pop-balance}
  (1 - \varepsilon)\bar p \;\le\; p(V_i) \;\le\; (1 + \varepsilon)\bar p,
  \qquad \bar p := \frac{p(V)}{k},
\end{equation}
for a tolerance $\varepsilon \in [0, 1)$. We write $\mathcal{P}_k(G; \varepsilon)$
for the set of valid plans. The \emph{cut edges} of a plan are
\[
  \partial P \;=\; \{(u,v) \in E : P(u) \ne P(v)\},
\]
and we call $|\partial P|$ the \emph{cut-edge count} of $P$. This
quantity is the discrete analog of the perimeter of the district
boundaries and serves as our primary compactness
metric~\cite{duchin2018discrete}.

\begin{figure}[H]
  \centering
  \includegraphics[width=0.5\textwidth]{fig/redistricting.png}
  \caption{The dual graph of Iowa's 99 counties (left) and the four
    congressional districts enacted after the 2010 census (right).
    Nodes sit at county centroids; two counties are joined by an edge
    when they share a boundary segment.}
  \label{fig:iowa-dual}
\end{figure}

\subsection{Recombination (ReCom)}\label{sec:background-recom}

Recombination~\cite{deford2021recombination} is a Markov chain on
$\mathcal{P}_k(G; \varepsilon)$ whose single step is given in Algorithm~\ref{alg:recom}.

\begin{algorithm}[H]
\caption{Vanilla ReCom step (Chain V)}\label{alg:recom}
\KwIn{Current plan $P$, dual graph $G$, balance tolerance $\varepsilon$.}
\KwOut{Next plan $P'$.}
Select two adjacent districts $V_i, V_j$ uniformly at random from boundary pairs;\\
Form the induced subgraph $H = G[V_i \cup V_j]$;\\
\Repeat{a cut $e^*$ exists yielding two $\varepsilon$-balanced components}{
  Sample a uniform random spanning tree $T$ of $H$ via Wilson's algorithm~\cite{wilson1996generating};\\
  Search edges of $T$ for a cut $e^*$ whose removal produces two $\varepsilon$-balanced components.
}
Define $V_i', V_j'$ as the two components of $T \setminus \{e^*\}$;\\
$P' \gets P$ with $V_i, V_j$ replaced by $V_i', V_j'$;\\
\Return $P'$.
\end{algorithm}

\noindent The tree-based proposal flips many precincts simultaneously
between $V_i$ and $V_j$ while preserving contiguity automatically (any
edge removal from a tree produces two connected components). This is
the key advantage over single-precinct flip chains, which traverse the
state space too slowly to be practically useful.

\textbf{Chains compared.} We compare six chains, three of them from
this work (Table~\ref{tab:five-chains}).

\begin{table}[H]
\centering
\small
\begin{tabular}{llll}
\toprule
Chain & Method & Mechanism & Stationarity claim \\
\midrule
V & Vanilla ReCom~\cite{deford2021recombination} & spanning-tree split, always accept & unknown target \\
A & Tempered ReCom & ReCom + MH on cuts (symmetric-$q$) & approximate \\
B & ReCom + QP-diffusion mix (this work) & continuous QP, Cheeger sweep, repair & none (heuristic) \\
C & Tempered ReCom + IQP local (this work) & MIQP refinement, orphan-flip repair & none (matheuristic) \\
D & ReCom + spectral mix (after~\cite{davies2025spectral}) & Fiedler + balanced sweep, no $\beta$, no repair & none (heuristic) \\
E & QP-MH (this work, \S\ref{sec:qp-mh}) & QP + Bernoulli rounding + exact MH & \textbf{exact} \\
\bottomrule
\end{tabular}
\caption{The six chains compared in this paper. The last column is the
  chain's relationship to the tempered target
  $\pi(P) \propto \exp(-\lambda|\partial P|)$. Only Chain E carries an
  exact Metropolis--Hastings correction; Chains B--D are heuristics
  whose output distributions we characterize empirically, and Chain A
  relies on the symmetric-$q$ approximation
  (\S\ref{sec:qp-mh}).}
\label{tab:five-chains}
\end{table}

Tempered ReCom uses a Metropolis--Hastings filter targeting
$\pi(P) \propto \exp(-\lambda |\partial P|)$ with the symmetric-$q$
approximation (treating the forward and reverse proposal probabilities
as equal); see Appendix~\ref{app:hyperparams} for the locked-in
$\lambda = 0.10$. Our Chain D departs from published SpecReCom in one
deliberate way. Instead of the sign split of the Fiedler vector (which
ignores population balance), it runs our balance-constrained Cheeger
sweep over the Fiedler order, closer in spirit to BalSpecReCom, so that all chains produce valid plans and the comparison isolates the
role of the fidelity term. All chains also enforce the global
two-sided balance window of Equation~\eqref{eq:balance-window} on
every move; we note (\S\ref{sec:exp-balance}) that the raw GerryChain
\texttt{recom} proposal alone does not.

\paragraph{The proposal we replace.}\label{sec:background-weakness}
The tree-and-cut step is kernel-blind (\S\ref{sec:intro-problem}),
and the cut edge is chosen only after the tree is drawn, so the
proposal cannot anticipate the shape of the split. The diffusion
proposal of \S\ref{sec:model} replaces that step with a single convex
optimization of a kernel-weighted Laplacian energy plus a
locality-preserving fidelity term. The energy
$x^\top L\, x = \sum_{(u,v) \in E} w_{uv} (x_u - x_v)^2$ equals the
weighted cut count for binary $x$, so minimizing it directly
minimizes the discrete perimeter in the chosen kernel.

%════════════════════════════════════════════════════════════════════════════════
\section{The QP-diffusion model: continuous and integer formulations}\label{sec:model}
%════════════════════════════════════════════════════════════════════════════════

\subsection{Locality kernels}\label{sec:kernels}

Choosing the edge weights $w_{ij}$ is the most important modeling
decision. The kernel determines what \emph{smoothness} means and
therefore controls how district boundaries move under the relaxation.
We restrict all weights to geographic neighbors,
\begin{equation}\label{eq:kernel-local}
  w_{ij} = 0 \quad \text{if } (v_i, v_j) \notin E,
\end{equation}
ensuring that diffusion operates only through geographic adjacency and
preventing long-range interactions that could produce disconnected
districts.

\paragraph{Uniform kernel ($k_{\mathrm{unif}}$).} $w_{ij} = 1$ for
every edge. The Laplacian $L$ is the combinatorial graph Laplacian;
$x^\top L x$ is exactly the cut edge count. In the $\beta \to 0^+$
limit the proposal approaches a balanced Fiedler-style split
(\S\ref{sec:iqp-subsumption}).

\paragraph{Perimeter kernel ($k_{\mathrm{perim}}$).} $w_{ij} = |\partial v_i \cap \partial v_j|$
(shared-boundary length). The Laplacian quadratic form $x^\top L x$
is then the \emph{geographic perimeter} of the partition, weighted by
boundary segment length. This kernel discourages fragmented districts
and favors boundaries that cut along short segments.

\paragraph{Density kernel ($k_{\mathrm{dens}}$).} Let $\rho_v$ denote
the population density at vertex $v$ (population per unit area). A
Gaussian density kernel
\[
  w_{ij} \;=\; |\partial v_i \cap \partial v_j| \cdot
  \exp\!\Bigl( -\frac{(\rho_i - \rho_j)^2}{\sigma^2} \Bigr)
\]
encourages cuts along sharp density gradients (urban--rural
boundaries) while preserving the topological role of shared boundary
length.

\paragraph{Demographic kernel ($k_{\mathrm{demog}}$).} Let
$d_v \in \R^m$ be a vector of demographic attributes. A Gaussian
similarity kernel $w_{ij} = \exp(-\|d_i - d_j\|^2 / \sigma^2)$
encourages communities of interest. We define this kernel for
completeness but do not use it in our experiments, and a caution
stronger than ``use as a soft regularizer'' is warranted. If the
attribute vector encodes race, a proposal that systematically shapes
district boundaries around racial similarity risks making race the
predominant districting criterion, which triggers strict scrutiny
under the \emph{Shaw v.\ Reno} line of racial-gerrymandering cases.
Any practical use of demographic kernels should be reviewed against
that doctrine, should document the kernel's weight relative to
race-neutral criteria, and is more defensible for attributes that are
legitimate districting factors in the jurisdiction (e.g., municipal
membership, communities of interest as defined by state law).

\paragraph{Composite kernels.} Linear combinations $w_{ij} = \sum_k a_k w_{ij}^{(k)}$
with $a_k \ge 0$ remain valid, since they preserve nonnegativity and locality.

\subsection{The convex QP}\label{sec:qp-model}

Let $H = G[V_i \cup V_j]$ be the merged subgraph of two adjacent
districts. Let $L_H$ be the kernel-weighted Laplacian of $H$, let
$x_0 \in \{0,1\}^{|V_H|}$ encode the current bipartition (1 for
vertices in $V_i$, 0 for vertices in $V_j$), and let $p$ be the
population vector restricted to $H$. Fix weights $\alpha, \beta > 0$
and balance tolerance $\varepsilon \in (0, 1)$. Crucially, the balance
constraint is anchored to the \emph{global} per-district ideal
$\bar p = p(V)/k$ from Equation~\eqref{eq:pop-balance}, not to
$p(V_H)/2$. Both sides of the split must individually satisfy the
plan-validity bound, which yields the two-sided window
\begin{equation}\label{eq:balance-window}
  \ell \;=\; \max\bigl((1-\varepsilon)\bar p,\; p(V_H) - (1+\varepsilon)\bar p\bigr),
  \qquad
  u \;=\; \min\bigl((1+\varepsilon)\bar p,\; p(V_H) - (1-\varepsilon)\bar p\bigr).
\end{equation}
Anchoring to $p(V_H)/2$ instead (a tempting simplification, and one
present in an earlier version of our own implementation) only bounds
the local split. Because $p(V_H)$ itself drifts as the chain moves
population between pairs, district populations can ratchet outside the
global tolerance by a factor of up to $(1+\varepsilon)$ per move on
heterogeneous graphs. Our experiments audit the global bound at every
step (\S\ref{sec:experiments}).

\begin{definition}[Continuous QP-diffusion]\label{def:qp}
The continuous QP-diffusion bipartition of $H$ is
\begin{align}
  x^\star \;=\; \arg\min_{x \in [0,1]^{|V_H|}} \;\;
    & \alpha\, x^\top L_H\, x \;+\; \beta\, \|x - x_0\|_2^2 \label{eq:qp-obj} \\
  \text{subject to} \;\;
    & \ell \;\le\; p^\top x \;\le\; u,
    \label{eq:qp-balance}
\end{align}
with $\ell, u$ as in Equation~\eqref{eq:balance-window}.
\end{definition}

The first term in~\eqref{eq:qp-obj} is the kernel-weighted Laplacian
smoothness energy. For binary $x$ it equals the weighted cut count
$\sum_{(u,v) \in E(H)} w_{uv}\,\mathbf{1}[x_u \ne x_v]$, and the
continuous relaxation $x \in [0,1]^{|V_H|}$ extends this to a smooth
objective on the unit cube. The second term is the \emph{fidelity term},
which penalizes deviation from the current bipartition and ensures
that each proposed move is local. The balance constraint enforces
$\varepsilon$-balance on the proposed split.

\begin{proposition}[Convexity and well-posedness]\label{prop:convex}
The objective in~\eqref{eq:qp-obj} is strictly convex on
$[0,1]^{|V_H|}$ (since $\beta > 0$ and $L_H$ is positive semidefinite).
The feasible region defined by~\eqref{eq:qp-balance} is a nonempty
polytope whenever the current plan is valid (it then contains $x_0$).
Hence the QP admits a unique minimizer $x^\star$, computable in
polynomial time by any convex QP solver.
\end{proposition}

\paragraph{Spectral interpretation.} Projecting the interior optimality
condition $(\alpha L_H + \beta I) x^\star = \beta x_0$ onto the
eigenbasis $\{\phi_k\}$ of $L_H$ with eigenvalues
$0 = \lambda_1 \le \lambda_2 \le \cdots \le \lambda_n$ yields the
closed-form frequency response
\begin{equation}\label{eq:freq-response}
  \hat x^\star_k \;=\; \frac{\beta}{\alpha \lambda_k + \beta}\,
  \hat x_{0,k},
  \qquad \hat y_k := \phi_k^\top y.
\end{equation}
The QP therefore acts as a low-pass spectral filter on the prior $x_0$,
with cutoff controlled by the ratio $\alpha / \beta$. Equivalently, with
$\tau := \alpha / \beta$, the optimum satisfies
$(x^\star - x_0)/\tau = -L_H x^\star$, the implicit Euler
discretization of the graph heat equation $\dot x = -L_H x$. Inspecting
the leading eigenvectors $\phi_2, \phi_3, \ldots$ provides a sanity
check on kernel choice. If they look like noise rather than coherent
geographic regions, the chosen kernel is probably not informative.

\subsection{Theory of the diffusion proposal}\label{sec:diffusion-theory}

The spectral identity is usually read as interpretation. Worked
through, it predicts how the proposal behaves, and every prediction
below is checked against the solver. Write $\tau := \alpha/\beta$ for
the diffusion timestep, $L = L_H$ for the kernel-weighted Laplacian of
the merged pair, $n = |V_H|$, and $\mathbf{1}_i \in \{0,1\}^{n}$ for
the indicator of $V_i$, so that $x_0 = \mathbf{1}_i$.

\begin{lemma}[The box constraints never bind]\label{lem:row-stochastic}
For any nonnegative kernel and any $\tau > 0$, the matrix
$(I + \tau L)^{-1}$ is symmetric, entrywise nonnegative, and has unit
row sums. Hence the unconstrained minimizer
$x^\star = (I + \tau L)^{-1} x_0$ of \eqref{eq:qp-obj} is a convex
combination of the entries of $x_0$ and lies in $[0,1]^{n}$, and the
QP of Definition~\ref{def:qp} reduces to this single linear solve
whenever the balance constraint \eqref{eq:qp-balance} is slack at
$x^\star$.
\end{lemma}

\begin{proof}
$I + \tau L$ is symmetric with positive diagonal, nonpositive
off-diagonal entries, and row sums equal to one, so it is strictly
diagonally dominant, hence a nonsingular M-matrix, and its inverse is
entrywise nonnegative. Since $L\mathbf{1} = 0$ we have
$(I + \tau L)\mathbf{1} = \mathbf{1}$ and so
$(I + \tau L)^{-1}\mathbf{1} = \mathbf{1}$. Each row of the inverse is
therefore a probability vector and $x^\star_v$ is a weighted average of
entries of $x_0 \in \{0,1\}^n$. The objective is strictly convex, so if
the unconstrained minimizer satisfies the one remaining constraint it
is the constrained minimizer.
\end{proof}

\begin{proposition}[Expected move size]\label{prop:move-size}
Under the single-shot Bernoulli rounding of \S\ref{sec:qp-mh}, and
whenever the balance constraint is slack, the expected number of units
of $H$ that a proposal reassigns is
\begin{equation}\label{eq:move-size}
  \E[\#\mathrm{flips}]
  \;=\; \|x^\star - x_0\|_1
  \;=\; 2\,\mathbf{1}_i^{\top}(I + \tau L)^{-1}\mathbf{1}_j
  \;=\; \frac{2\,|V_i|\,|V_j|}{n}
        \;-\; 2\sum_{k \ge 2}\frac{(\phi_k^{\top}\mathbf{1}_i)^2}{1 + \tau\lambda_k},
\end{equation}
where $(\lambda_k, \phi_k)$ are the eigenpairs of $L$ with
$\lambda_1 = 0$.
\end{proposition}

\begin{proof}
A unit with $x_{0,v} = 0$ flips with probability $x^\star_v$ and one
with $x_{0,v} = 1$ with probability $1 - x^\star_v$; both equal
$|x^\star_v - x_{0,v}|$ because $x^\star_v \in [0,1]$ by
Lemma~\ref{lem:row-stochastic}, which gives the first equality.
Write $P = (I + \tau L)^{-1}$. For $v \in V_i$,
$1 - x^\star_v = \sum_{u \in V_j} P_{vu}$ by unit row sums, and for
$v \in V_j$, $x^\star_v = \sum_{u \in V_i} P_{vu}$; summing and using
symmetry of $P$ gives $2\,\mathbf{1}_i^{\top} P\, \mathbf{1}_j$.
Expanding $P = \sum_k \phi_k\phi_k^{\top}/(1 + \tau\lambda_k)$ with
$\phi_1 = \mathbf{1}/\sqrt{n}$ gives the $k = 1$ term
$|V_i||V_j|/n$, and since $\phi_k \perp \mathbf{1}$ for $k \ge 2$,
$\phi_k^{\top}\mathbf{1}_j = -\phi_k^{\top}\mathbf{1}_i$, which
gives the sum. Clipping at $\delta = 10^{-12}$ changes the expectation
by at most $\delta n$.
\end{proof}

\begin{corollary}[Small-timestep law, monotonicity, and ceiling]\label{cor:small-tau}
Let $|\partial_w(V_i, V_j)|$ be the boundary of the pair measured in
the proposal's kernel. Then $\E[\#\mathrm{flips}]$ is increasing in
$\tau$,
\[
  \E[\#\mathrm{flips}] \;=\; 2\tau\,|\partial_w(V_i,V_j)| + O(\tau^2)
  \quad (\tau \to 0),
  \qquad
  \E[\#\mathrm{flips}] \;\le\; \frac{2\,|V_i|\,|V_j|}{n}
  \quad\text{for every } \tau,
\]
with equality in the limit $\tau \to \infty$. For a balanced pair the
ceiling is $n/2$, half the merged pair.
\end{corollary}

\begin{proof}
Each subtracted term in \eqref{eq:move-size} is nonnegative and
decreasing in $\tau$, which gives monotonicity and the bound, and the
terms vanish as $\tau \to \infty$. For the small-$\tau$ law, expand
$(1 + \tau\lambda_k)^{-1} = 1 - \tau\lambda_k + O(\tau^2)$; the
$\tau^0$ terms cancel because
$\sum_{k \ge 2}(\phi_k^{\top}\mathbf{1}_i)^2 = |V_i| - |V_i|^2/n =
|V_i||V_j|/n$, and the $\tau^1$ term is
$2\tau\sum_k \lambda_k(\phi_k^{\top}\mathbf{1}_i)^2 =
2\tau\,\mathbf{1}_i^{\top}L\,\mathbf{1}_i = 2\tau|\partial_w|$.
\end{proof}

Three consequences follow. The continuous proposal is far cheaper
than ``solve a QP'' suggests, since it is one linear solve in the
common case. The move size is a design dial with a known scale, so to
propose moves of about $m$ units one sets
$\alpha/\beta \approx m / (2|\partial_w|)$ while the small-timestep law
holds. And the dial has a hard stop. No choice of $\alpha$ and
$\beta$ can lift the expected move above half the merged pair, so
mobility cannot be bought by tuning, which is the structural reason a
global move is needed alongside the exact local chain
(\S\ref{sec:mobility}, \S\ref{sec:disc-future}). The kernel enters
through $|\partial_w|$, so a perimeter kernel takes larger steps
exactly where the geographic border is long.

\begin{figure}[H]
  \centering
  \includegraphics[width=\textwidth]{fig/spectral_move_size.pdf}
  \caption{The move-size law against the solver. (a)~On the enumerated
    $5 \times 5$ fiber, 1\,411 merged pairs drawn from 200 plans; the
    exact filter norm of \eqref{eq:move-size} and the small-timestep
    law $2\tau|\partial_w|$ agree for $\tau \le 0.03$ and separate as
    the gain saturates. Raising $\tau$ by $1000\times$ raises the move
    by $76\times$, and the last decade buys $2.1\times$. The locked
    Chain E setting $\tau = 0.1$ is marked. (b)~The same on Iowa, 127
    pairs from 25 tree-partition plans. (c)~Fraction of merged pairs
    whose balance constraint is active. On the uniform grid it never
    is, so the proposal is a linear solve at every step; on real county
    geography it is active for 41\% of pairs at $\tau = 0.1$ and 86\%
    at $\tau = 10$, so there the population constraint, not the
    diffusion, caps the move.}
  \label{fig:spectral-move-size}
\end{figure}

\textbf{Verification.} Across 10\,766 linear solves on the two test
beds the unconstrained optimum never left the unit cube, as
Lemma~\ref{lem:row-stochastic} requires. On 120 merged pairs with
balance slack, Gurobi's constrained solution and the linear solve
agree to $1.1 \times 10^{-4}$, the solver tolerance, so the
implementation matches the theory. The median relative error of the
small-timestep law against the exact filter norm is $2.5\%$ at
$\tau = 0.01$, $7.4\%$ at $0.03$, and $24\%$ at $0.1$
(Figure~\ref{fig:spectral-move-size}). At the paper's locked Chain E
setting the expected move is $0.42$ units on the $5 \times 5$ fiber
and $0.82$ on Iowa, under one unit per proposal, which is the
quantitative content of the chain's stickiness reported in
\S\ref{sec:exp-5x5}.

\subsection{Cheeger-sweep rounding and orphan-flip repair}\label{sec:rounding}

The continuous optimum $x^\star \in [0,1]^{|V_H|}$ must be rounded to a
binary bipartition. We use a \emph{Cheeger sweep}~\cite{chung1997spectral}. Sort the entries of $x^\star$ in descending order, and for each
threshold $k = 1, \ldots, |V_H| - 1$ form the candidate set $S_k$ of the
top-$k$ vertices. Among candidates that satisfy $\varepsilon$-balance,
return the one minimizing the cut count
$|\partial(S_k, V_H \setminus S_k)|$. Candidates with disconnected sides
are passed through the orphan-flip repair below before evaluation.

\begin{algorithm}[H]
\caption{Cheeger-sweep rounding with orphan-flip repair}\label{alg:cheeger-sweep}
\KwIn{Continuous optimum $x^\star \in [0,1]^{|V_H|}$, population
  vector $p$, target $\bar p$, tolerance $\varepsilon$, merged subgraph $H$.}
\KwOut{Binary bipartition $S \subset V_H$ or \texttt{FAIL}.}
Sort $V_H$ by $x^\star$ descending; call the order $v_{(1)}, \ldots, v_{(|V_H|)}$;\\
$\mathcal{C} \gets \emptyset$ \tcp*{candidates}
\For{$k = 1$ \KwTo $|V_H| - 1$}{
  $S_k \gets \{v_{(1)}, \ldots, v_{(k)}\}$;\\
  \uIf{$p(S_k) \in [(1-\varepsilon)\bar p, (1+\varepsilon)\bar p]$ and both sides connected}{
    Add $S_k$ to $\mathcal{C}$ with cut $|\partial(S_k, V_H \setminus S_k)|$;
  }\ElseIf{$p(S_k)$ in balance and sides have $> 1$ components}{
    $S_k' \gets$ orphan-flip repair of $S_k$ (Algorithm~\ref{alg:orphan-flip});\\
    \If{$S_k'$ satisfies balance and connectivity}{
      Add $S_k'$ to $\mathcal{C}$;
    }
  }
}
\lIf{$\mathcal{C} = \emptyset$}{\Return \texttt{FAIL}}
\Return $\arg\min_{S \in \mathcal{C}} |\partial(S, V_H \setminus S)|$.
\end{algorithm}

The sweep can return a candidate $S_k$ whose induced subgraph $G[S_k]$
or $G[V_H \setminus S_k]$ is disconnected. A single boundary vertex
on the wrong side of the threshold is enough. Discarding such
candidates outright would inflate the rejection rate; instead we repair
them. The \emph{orphan-flip} routine identifies the disconnected
components of each side, keeps the largest component, and flips every
smaller (\emph{orphan}) component to the opposite side. Population
balance is then re-verified; if the repair preserves both balance and
connectivity, the repaired candidate enters the Cheeger sweep's
candidate pool. Algorithm~\ref{alg:orphan-flip} formalizes this procedure.

\begin{algorithm}[H]
\caption{Orphan-flip repair}\label{alg:orphan-flip}
\KwIn{Candidate bipartition $(S, V_H \setminus S)$.}
\KwOut{Repaired bipartition or \texttt{FAIL}.}
Identify connected components of $G[S]$ and $G[V_H \setminus S]$;\\
\For{each side}{
  Keep the largest connected component;\\
  Flip every smaller component to the opposite side;
}
Verify both final sides are connected and $\varepsilon$-balanced;\\
\lIf{verification fails}{\Return \texttt{FAIL}}
\Return repaired $(S', V_H \setminus S')$.
\end{algorithm}

\textbf{Why repair matters.} On heterogeneous-population graphs such
as Iowa, the unrepaired Cheeger sweep produces near-feasible
candidates whose few orphan vertices (often singletons at the
boundary) can be reassigned without affecting the cut count.
Without repair, post-rounding rejection rates on Iowa run at 5--15\%;
with repair, $<3\%$. SpecReCom~\cite{davies2025spectral} does not use
repair (and we keep it out of our faithful Chain D
implementation); without the $\beta$ fidelity term, repair tends to
accept lower-quality cuts.

\paragraph{On sweep guarantees.} We do not claim the Cheeger
bound~\cite{chung1997spectral} for the rounded solution. The sweep is
applied to the constrained proximal optimum of
Definition~\ref{def:qp}, not to the Fiedler vector, and the balance
restriction on admissible thresholds voids the standard argument, so
the sweep is a rounding heuristic whose output we measure. It also
selects the candidate minimizing the \emph{unweighted} cut even when
the energy uses a non-uniform kernel, a configuration choice we flag
in \S\ref{sec:exp-multikernel}.

\subsection{Hybrid composition with ReCom (Chain B)}\label{sec:hybrid-mh}

Chain B's single step is a random mixture. With probability
$p_R = 0.5$, take a Vanilla ReCom step (Algorithm~\ref{alg:recom});
otherwise take a QP-diffusion step (Definition~\ref{def:qp} and
Algorithm~\ref{alg:cheeger-sweep}). We state plainly what Chain B is
and is not. \emph{Neither branch carries a Metropolis correction}, the ReCom branch always accepts (as in Vanilla ReCom), and the
QP-diffusion branch applies the rounded, repaired split
unconditionally. Chain B therefore has no stationarity claim beyond
ergodicity; it is an exploratory heuristic whose empirical output
distribution we characterize against enumerated ground truth in
\S\ref{sec:exp-5x5}. The exactly corrected sampler built from the same
proposal is Chain E (\S\ref{sec:qp-mh}).

\begin{proposition}[Ergodicity of the hybrid]\label{prop:ergodic}
If ReCom is ergodic on $\mathcal{P}_k(G; \varepsilon)$ (which holds
under standard assumptions; see~\cite{deford2021recombination,
cannon2024irreducibility}), then the mixture of ReCom and
the QP-diffusion step with any $p_R \in (0, 1)$ is also ergodic.
\end{proposition}

\begin{proof}
Write the mixture kernel as $K = p_R K_{\mathrm{R}} +
(1 - p_R) K_{\mathrm{Q}}$. If $K_{\mathrm{R}}^{T}(P, P') > 0$ for
some $T$, then $K^{T}(P, P') \ge p_R^{T}\, K_{\mathrm{R}}^{T}(P, P')
> 0$, so irreducibility of ReCom implies irreducibility of the
mixture. For aperiodicity, the ReCom branch already holds in place
with positive probability. The current split of any selected pair is
a balanced bipartition of the connected merged subgraph $H$, and it
is produced whenever the sampled spanning tree consists of a spanning
tree of each side joined by one crossing edge, an event of positive
probability under the uniform tree measure. Hence
$K_{\mathrm{R}}(P, P) > 0$, so $K(P, P) > 0$ for every $P$ and the
chain is aperiodic.
\end{proof}

\subsection{The exactly reversible QP-MH chain (Chain E)}\label{sec:qp-mh}

The deterministic rounding of \S\ref{sec:rounding} is what blocks a
Metropolis correction. An argmin-plus-sweep map has zero reverse
proposal probability, so the Hastings ratio degenerates. Replacing the
sweep with \emph{single-shot randomized rounding} restores a tractable,
everywhere-positive proposal density and yields an exactly reversible
chain. One step of Chain E is as follows.

\begin{enumerate}
  \item Pick an adjacent district pair $(V_i, V_j)$ uniformly at random
    from the $|\mathrm{bp}(P)|$ boundary pairs of the current plan $P$;
  \item Solve the (deterministic) QP of Definition~\ref{def:qp} on
    $H$ for $x^\star \in [0,1]^{|V_H|}$;
  \item Draw the proposal $x' \in \{0,1\}^{|V_H|}$ by independent
    Bernoulli rounding from the \emph{clipped} optimum,
    $x'_v \sim \mathrm{Bern}\bigl(\mathrm{clip}(x^\star_v;
    \delta, 1-\delta)\bigr)$ with $\delta = 10^{-12}$,
    \emph{once}, with no retry on infeasibility. Clipping is part of the
    proposal's definition. The QP optimum can have active box bounds
    (coordinates exactly $0$ or $1$, e.g.\ at interior vertices pinned
    by the fidelity term), and sampling from the same clipped vector
    whose density enters the acceptance ratio is what makes the
    proposal law and its stated density identical, with every binary
    $x'$ receiving positive probability;
  \item Accept with probability
    \begin{equation}\label{eq:qpmh-accept}
      \min\Bigl(1,\;
      e^{-\lambda(|\partial P'| - |\partial P|)}\,
      \frac{q(x_0 \mid P', \{i,j\})}{q(x' \mid P, \{i,j\})}\,
      \frac{|\mathrm{bp}(P)|}{|\mathrm{bp}(P')|}\Bigr),
    \end{equation}
    where $q(\cdot \mid P, \{i,j\})$ is the Bernoulli product density
    induced by the QP optimum at the corresponding state, and proposals
    that violate contiguity or the balance
    window~\eqref{eq:balance-window} are rejected outright.
\end{enumerate}

\begin{proposition}[Detailed balance of Chain E]\label{prop:qpmh-db}
Chain E satisfies detailed balance with respect to
$\pi(P) \propto \exp(-\lambda |\partial P|)$ restricted to
$\mathcal{P}_k(G; \varepsilon)$.
\end{proposition}

Throughout, fix the deterministic edge kernel $w$ (weight jitter is
never used in this chain) and write $\Omega = \mathcal{P}_k(G;
\varepsilon)$ for the space of valid plans. Plans are \emph{labeled}
assignments in this subsection, so that a Bernoulli draw and its
complement are distinct states and the proposal map is a bijection;
on unlabeled partitions the map is two-to-one and the transition
probability would need both terms. Since $\pi$ depends only on
$|\partial P|$ it is label-invariant, and forgetting labels after
sampling recovers the unlabeled law. For a plan $P$ and an adjacent
pair $\{i,j\} \in \mathrm{bp}(P)$ write $H = V_i \cup V_j$ and let
$x_0(P) \in \{0,1\}^{V_H}$ denote the indicator of $V_i$ on $H$ (the
pair is drawn from a sorted list, so which district plays the role of
$i$ is a deterministic function of the unordered pair). Say that
$P' \in \Omega$ is \emph{pair-adjacent} to $P \in \Omega$ if
$P' \ne P$, the two plans differ on exactly two districts $\{i,j\}$,
and $H = V_i \cup V_j$ is connected. This relation is clearly
symmetric in $P$ and $P'$.

\begin{lemma}[Proposal regularity]\label{lem:qpmh-regular}
Let $P, P' \in \Omega$ be pair-adjacent via districts $\{i,j\}$, and
suppose $\beta > 0$. Then:
\begin{enumerate}
  \item[(i)] $\{i,j\}$ is the unique unordered pair whose selection
    can generate the move $P \to P'$ or the move $P' \to P$, and it
    is selectable from both endpoints:
    $\{i,j\} \in \mathrm{bp}(P) \cap \mathrm{bp}(P')$.
  \item[(ii)] The QP of Definition~\ref{def:qp} on $H$ with the
    global window~\eqref{eq:balance-window} has the same (nonempty)
    feasible set for the forward and reverse solves, and its optimum
    is unique, hence a deterministic function of
    $(\text{endpoint}, \{i,j\})$.
  \item[(iii)] The clipped-Bernoulli density
    $q(\cdot \mid \cdot, \{i,j\})$ is bounded below by
    $\delta^{|V_H|} > 0$ on every binary labeling of $V_H$; in
    particular $q(x_0(P') \mid P, \{i,j\}) > 0$ and
    $q(x_0(P) \mid P', \{i,j\}) > 0$.
\end{enumerate}
\end{lemma}

\begin{proof}
(i) A step that selects pair $\{a,b\}$ relabels vertices only
between districts $a$ and $b$ and changes no other district. If
$P' \ne P$, at least one vertex of $H$ switched between $i$ and $j$,
so districts $i$ and $j$ both changed; any generating pair must
therefore equal $\{i,j\}$. For selectability, note that
$(V_i, V_j)$ and $(V_i', V_j')$ are both bipartitions of the
connected graph $H$ into nonempty parts (all four sets are nonempty
districts of valid plans), and every such bipartition is crossed by
at least one edge of $H$; hence the two districts are adjacent in
both plans.

(ii) The feasible set is $\{x \in [0,1]^{V_H} : \ell \le p^\top x
\le u\}$ with $\ell, u$ given by~\eqref{eq:balance-window}. It
depends only on $(H, \varepsilon, \bar{p})$, not on the endpoint, so
the forward and reverse solves optimize over the same set. It is
nonempty at both endpoints because the window
of~\eqref{eq:balance-window} is precisely the condition that both
sides of the split be valid districts, which $x_0(P)$ and $x_0(P')$
satisfy by validity of $P$ and $P'$. This is where globality is
essential. Anchoring the window to the merged-pair half
$(1 \pm \varepsilon)\,p(V_H)/2$ instead lets a valid plan fail its
own pair's window, zeroing a reverse density whose forward
counterpart is positive and destroying reversibility outright. The
objective $\alpha\, x^\top L_H x + \beta\, \|x - x_0\|^2$ has Hessian
$2(\alpha L_H + \beta I) \succ 0$ for $\beta > 0$, so the minimizer
over the (convex, compact, nonempty) feasible set exists and is
unique, and it is a function of $(L_H, x_0, \ell, u)$ only. Since
the kernel $w$ is deterministic, $L_H$ depends only on $H$, and
$x_0$ only on the endpoint and the pair.

(iii) Clipping maps the optimum into $[\delta, 1-\delta]^{V_H}$, so
every Bernoulli factor of the product density lies in
$[\delta, 1-\delta]$ and every outcome in $\{0,1\}^{V_H}$ has
probability at least $\delta^{|V_H|}$.
\end{proof}

\begin{proof}[Proof of Proposition~\ref{prop:qpmh-db}]
Let $K$ denote the transition kernel of Chain E and fix distinct
$P, P' \in \Omega$. If $P$ and $P'$ are not pair-adjacent, no single
step can produce either move. A step changes at most the two selected
districts, so a generating pair would have to be the differing pair;
if the plans differ on more than two districts no pair qualifies, and
if the union $H$ of the two differing districts is disconnected, that
pair is not adjacent in either plan and is never selected. Hence
$K(P, P') = K(P', P) = 0$ and detailed balance holds trivially.

Otherwise let $\{i,j\}$ be the unique generating pair, abbreviate
$q_{\to} = q(x_0(P') \mid P, \{i,j\})$ and
$q_{\gets} = q(x_0(P) \mid P', \{i,j\})$, and note that both are
positive and finite by Lemma~\ref{lem:qpmh-regular}(ii)--(iii).
Because the rounding is single-shot, $q_{\to}$ is the exact
probability that the proposal from $P$ with pair $\{i,j\}$ equals
$P'$. No renormalization by a feasibility event is involved (a
retry-until-feasible loop would replace $q$ by a conditional law
whose state-dependent normalizing constant does not cancel in the
ratio), and clipping \emph{before} sampling makes the sampled law and
the density in~\eqref{eq:qpmh-accept} the same object. Hence
\[
  K(P, P') \;=\; \frac{1}{|\mathrm{bp}(P)|}\; q_{\to}\;
  \alpha(P, P'), \qquad
  \alpha(P, P') = \min\Bigl(1,\;
  \frac{\pi(P')}{\pi(P)}\,
  \frac{q_{\gets}}{q_{\to}}\,
  \frac{|\mathrm{bp}(P)|}{|\mathrm{bp}(P')|}\Bigr),
\]
using $\pi(P')/\pi(P) = e^{-\lambda(|\partial P'| - |\partial P|)}$.
Multiplying through,
\[
  \pi(P)\, K(P, P') \;=\;
  \min\Bigl(
  \frac{\pi(P)\, q_{\to}}{|\mathrm{bp}(P)|},\;
  \frac{\pi(P')\, q_{\gets}}{|\mathrm{bp}(P')|}
  \Bigr),
\]
an expression symmetric under exchanging $(P, q_{\to})$ with
$(P', q_{\gets})$. The identical computation started from $P'$ gives
the same right-hand side, so
$\pi(P) K(P, P') = \pi(P') K(P', P)$.

It remains to check that the rejection pathways do not disturb this
identity. A Bernoulli draw whose split leaves a side empty,
disconnected, or outside the window~\eqref{eq:balance-window} is a
proposal onto a state of zero target mass and is rejected with
probability one, contributing only to the diagonal $K(P, P)$, for
which detailed balance is vacuous; the same applies to in-space
proposals rejected by~\eqref{eq:qpmh-accept}. Rejection on solver
failure is symmetric because, by
Lemma~\ref{lem:qpmh-regular}(ii), the solve is a deterministic
function of $(H, \{i,j\}, \text{endpoint})$, so a failure that
suppresses the move $P \to P'$ suppresses the reverse move
identically. Detailed balance of $K$ with respect to $\pi$ on
$\Omega$ follows, and with it stationarity of
$\pi$~\cite{tierney1998note}.
\end{proof}

\begin{remark}[Scope, verification, and ergodicity]\label{rem:qpmh-scope}
The proof treats the QP solve as exact; in floating point the
minimizer is computed to solver tolerance, and determinism of the
implementation (fixed kernel, no jitter, sorted pair enumeration) is
what the argument actually uses. As an end-to-end check, we assembled
the full transition kernel of the implementation on an exhaustively
enumerated $2 \times 3$ grid ($k = 2$, tolerance $\varepsilon = 1/3$,
so that the 18 labeled valid states include $4{+}2$ as well as
$3{+}3$ splits) and verified
$\pi(P)K(P,P') = \pi(P')K(P',P)$ entrywise, with maximum relative
violation $1.5 \times 10^{-15}$ (machine precision;
\texttt{db\_check\_2x3.py}). Detailed balance
gives stationarity, not ergodicity. Within a selected pair every
binary relabeling of $H$ has positive proposal probability by
Lemma~\ref{lem:qpmh-regular}(iii), but irreducibility of Chain E
across pairs on general $(G, k, \varepsilon)$ is not established, so
convergence statements are for the chain restricted to its
communicating class. On the enumerated $5 \times 5$ fiber the chain's cut distribution is
indistinguishable from its target (\S\ref{sec:exp-longrun}), on the
$2 \times 3$ fiber a million-step run's full occupancy distribution
is (\S\ref{sec:exp-2x3}), and the principled route to global mobility
is the $\pi$-Metropolized mixture described below.
\end{remark}

Chain E trades mobility for exactness. Because a proposal must be
feasible in a single Bernoulli draw, and because
Corollary~\ref{cor:small-tau} caps how far a draw can reach, the
chain behaves as a strictly local sampler. \S\ref{sec:mobility}
measures exactly why, and the answer is feasibility rather than the
Metropolis ratio. We therefore present Chain E as the
\emph{correctness anchor} of the framework (the member of the family
whose stationarity is provable) rather than as a practical standalone
sampler. The principled route to mobility is precise. Mixing E with a
global kernel preserves $\pi$ only if that kernel is \emph{also}
$\pi$-invariant. RevReCom as published is reversible with respect to
its spanning-tree measure, not our tempered $\pi$, so a naive
E/RevReCom mixture would preserve neither; the correct construction
Metropolizes the tree move toward $\pi$ using its computable Hastings
ratio~\cite{cannon2022spanning,autry2023metropolized} and mixes
\emph{that} kernel with E. This is the natural next step
(\S\ref{sec:disc-future}).

\paragraph{On the tempered baseline (Chain A).} Our Tempered ReCom
baseline applies the Metropolis filter
$\min(1, \exp(-\lambda \Delta|\partial P|))$ over the ReCom proposal,
i.e., it \emph{assumes a symmetric proposal}. ReCom's proposal is not
symmetric; the true Hastings ratio involves spanning-tree counts of
the affected districts~\cite{autry2023metropolized,cannon2022spanning}, so Chain A targets $\pi$ only approximately. On the enumerated
$5 \times 5$ fiber this bias is visible and quantifiable
(\S\ref{sec:exp-5x5}; empirical mean cuts $\approx 16.8$ against the
exact target mean $18.56$); it is inherited from the spanning-tree
proposal's own compactness preference, and correcting it requires the
tree-count acceptance ratios of Metropolized forest
recombination~\cite{autry2023metropolized}. We keep the symmetric-$q$
baseline because it is what practitioners run, and we report its bias
rather than hiding it.

\subsection{Mobility of the exact chain: feasibility, not acceptance}\label{sec:mobility}

Proposition~\ref{prop:move-size} says how far a proposal reaches.
Mobility is that reach times the probability the proposal is kept,
and the second factor has two causes of different kinds. A Bernoulli
draw can produce a split that is disconnected or outside the balance
window, in which case it is rejected with probability one before any
coin is flipped. Or it can be a valid plan that loses the Metropolis
coin. We measured both on the enumerated $5 \times 5$ fiber, with
$\alpha = \tau$ and $\beta = 1$, over 3\,000 steps per timestep from
the same initial plan (Table~\ref{tab:mobility}).

\begin{table}[H]
\centering
\small
\begin{tabular}{rrrrrrrr}
\toprule
$\tau$ & $\E[\#\mathrm{flips}]$ & single-flip share & feasible $\mid$ single
  & $\Pr[\text{feasible}]$ & accept $\mid$ feasible & units moved / step & states \\
\midrule
0.01 & 0.036 & 0.035 & 0.000 & 0.964 & 1.000 & 0.0000 & 1 \\
0.03 & 0.126 & 0.110 & 0.000 & 0.884 & 1.000 & 0.0033 & 6 \\
0.10 & 0.385 & 0.268 & 0.000 & 0.683 & 0.999 & 0.0093 & 12 \\
0.30 & 0.909 & 0.349 & 0.000 & 0.437 & 0.989 & \textbf{0.0353} & \textbf{38} \\
1.00 & 1.914 & 0.254 & 0.000 & 0.193 & 0.983 & 0.0313 & 37 \\
3.00 & 2.672 & 0.168 & 0.000 & 0.074 & 0.986 & 0.0187 & 19 \\
\bottomrule
\end{tabular}
\caption{Chain E against the diffusion timestep on the $5 \times 5$
  fiber ($k = 5$, exact balance). ``Feasible $\mid$ single'' is the
  feasibility rate of single-unit proposals, ``accept $\mid$
  feasible'' the Metropolis acceptance rate among feasible
  proposals, and ``units moved / step'' the product of acceptance and
  accepted move size, the quantity a tuning rule should maximize.
  ``States'' counts distinct fiber states reached in 3\,000 steps.
  The locked setting of Appendix~\ref{app:hyperparams} is
  $\tau = 0.10$.}\label{tab:mobility}
\end{table}

Three facts follow. First, every single-unit proposal is infeasible at
every timestep, exactly as exact balance requires, since reassigning
one cell produces districts of sizes 4 and 6. The leading-order move
of the diffusion proposal is therefore always rejected on this fiber,
and the chain moves only through second-order coordinated draws. This
is the mechanism behind the coverage figures of \S\ref{sec:exp-5x5},
derived rather than observed. Second, the Metropolis ratio is never
the bottleneck. Given a feasible proposal, acceptance is between 98
and 100 percent throughout. The chain's stickiness is a feasibility
effect, and the natural phrase ``large moves are rarely accepted'' is
wrong in an instructive way. They are rarely feasible, and when
feasible they are accepted. Third, acceptance rate is a misleading
diagnostic for this chain. At $\tau = 0.01$ it is 96 percent while
the chain never leaves its initial state, because 96 percent of
proposals flip nothing. We report units moved per step instead.

Units moved per step has an interior maximum near $\tau = 0.3$. The
setting locked in Appendix~\ref{app:hyperparams}, $\tau = 0.1$, sits
about $3.8\times$ below it, so the chain as run in
\S\ref{sec:experiments} is mistuned by its own theory and can be
improved for free. Even at the optimum a thousand steps move about 35
units, and Corollary~\ref{cor:small-tau} says no timestep can lift the
per-proposal reach above half the merged pair. Both statements
matter. Chain E can be several times more mobile than shown, and
several times is not enough, which is why a global move is required
(\S\ref{sec:disc-future}). Under a loose tolerance single flips become
feasible and mobility gains a first-order term in $\tau$, which is why
the same chain climbs steadily on Iowa (Figure~\ref{fig:iowa-trace})
while freezing on the exact-balance grid.

Since acceptance given feasibility is close to one, mobility is
essentially $\sum_k k\,\Pr[k \text{ flips}]\,\Pr[\text{feasible} \mid
k]$, with a Poisson-binomial flip count whose parameters
Proposition~\ref{prop:move-size} supplies; a closed-form mobility law
is within reach and is left to future work.

\subsection{The integer QP-diffusion (Chain C)}\label{sec:iqp}

The integer counterpart of Definition~\ref{def:qp} replaces
$x \in [0,1]^{|V_H|}$ with $x \in \{0, 1\}^{|V_H|}$ and solves the
resulting mixed-integer convex quadratic program (MIQP) directly.

\paragraph{Definition and solver.}\label{sec:iqp-def}

\begin{definition}[Integer QP-diffusion]\label{def:iqp}
The IQP local refinement of $H$ is
\begin{align}
  x^\star \;=\; \arg\min_{x \in \{0,1\}^{|V_H|}} \;\;
    & \alpha\, x^\top L_H\, x \;+\; \beta\, \|x - x_0\|_2^2 \label{eq:iqp-obj} \\
  \text{subject to} \;\;
    & \ell \;\le\; p^\top x \;\le\; u,
\end{align}
with the same global-ideal balance window $(\ell, u)$ of
Equation~\eqref{eq:balance-window} as the continuous case.
\end{definition}

The objective is convex quadratic and the variable is binary.
Gurobi's MIQP solver handles this directly via branch-and-bound; in
practice, a 5-second time limit per solve is rarely hit on Iowa-sized
merged pairs ($\sim 30$--$60$ precincts).

\paragraph{Subsumption of cut-min; relation to SpecReCom.}\label{sec:iqp-subsumption}

\begin{proposition}[Special cases]\label{prop:iqp-special-cases}
Setting $\alpha = 1$, $\beta = 0$, $w_{ij} \equiv 1$ in
Definition~\ref{def:iqp} gives the classical minimum-cut integer
program subject to the balance window. Hence Chain C subsumes
the discrete cut-minimization MIP as a special case.
\end{proposition}

We are careful \emph{not} to claim the analogous statement for the
continuous case. At $\beta = 0$ the relaxation of
Definition~\ref{def:qp} is degenerate. Any feasible constant vector
$c\mathbf{1}$ attains energy zero, so the minimizer carries no cut
information and the sweep is undefined. SpecReCom's Fiedler direction
arises from the Rayleigh quotient \emph{restricted to}
$x \perp \mathbf{1}$, a constraint our proximal program does not
impose; it is the fidelity term ($\beta > 0$) that breaks the constant
degeneracy here, with $x_0$ playing the symmetry-breaking role that
orthogonality plays in spectral partitioning. The precise relationship
is that as $\beta \to 0^+$ with $x_0$ non-constant, the minimizer's
projection off $\mathbf{1}$ aligns with the low-frequency eigenspace
of $L_H$ (Equation~\eqref{eq:freq-response} with
$\beta/(\alpha\lambda_k + \beta) \to 0$ fastest at high frequencies),
so Chain B approaches a Fiedler-like split \emph{in direction} while
retaining balance feasibility. This gives a BalSpecReCom-style variant rather
than SpecReCom itself. Chain D implements that variant directly.

\paragraph{Locality via the fidelity term.}\label{sec:iqp-locality}

For $\beta = 0$, the IQP optimum is a minimum cut of $H$ subject only
to population balance and may differ arbitrarily from $x_0$. For
$\beta > 0$ the optimum is biased toward $x_0$, since every flip
incurs a $\beta$ penalty, so the proposed split stays near the current
one, is mostly contiguous already because $x_0$ is, and needs only a
few boundary units repaired; the rejection figures of
\S\ref{sec:rounding} are the measured consequence.

\paragraph{One step of Chain C.}\label{sec:chain-c-step}
A single step of Chain C wraps an outer Tempered ReCom move
(Algorithm~\ref{alg:recom} plus the Metropolis--Hastings filter of
\S\ref{sec:hybrid-mh}) around $n_{\mathrm{iqp}}$ deterministic
IQP local-refinement passes. Each inner pass picks an adjacent district
pair, solves the IQP on the merged subgraph, and applies orphan-flip
repair; if the repair fails the inner pass leaves the plan unchanged.
The full step is given as Algorithm~\ref{alg:chain-c}.

\begin{algorithm}[H]
\caption{One step of Chain C}\label{alg:chain-c}
\KwIn{Current plan $P$, parameters $\alpha, \beta, n_{\mathrm{iqp}}, \lambda, \varepsilon$.}
\KwOut{Next plan $P'$.}
$P_0 \gets $ Tempered ReCom step (Chain A) applied to $P$\;
$P_{\mathrm{curr}} \gets P_0$\;
\For{$t = 1$ \KwTo $n_{\mathrm{iqp}}$}{
  Select adjacent district pair $(V_i, V_j)$ uniformly from boundary pairs of $P_{\mathrm{curr}}$\;
  Form $H = G[V_i \cup V_j]$, build kernel-weighted $L_H$\;
  Solve IQP (Definition~\ref{def:iqp}) for $x^\star$\;
  Apply orphan-flip repair (Algorithm~\ref{alg:orphan-flip})\;
  \uIf{repaired bipartition feasible and connected}{
    Update $V_i, V_j$ in $P_{\mathrm{curr}}$ to the new bipartition\;
  }\Else{
    $P_{\mathrm{curr}}$ unchanged\;
  }
}
\Return $P' \gets P_{\mathrm{curr}}$.
\end{algorithm}

\paragraph{What Chain C is, and is not.}\label{sec:chain-c-honest}
The IQP loop is \emph{deterministic} given the random pair choices. There is no Metropolis acceptance inside the loop, and the outer
acceptance applies only to the Tempered ReCom move. We are explicit
about the consequence. \textbf{Chain C is not a sampler for the
tempered target} $\pi(P) \propto \exp(-\lambda|\partial P|)$, and no
completion of the construction can make it one. A deterministic map
$T$ composed after a $\pi$-invariant kernel preserves $\pi$ only if
$T$ is itself $\pi$-preserving (essentially a $\pi$-preserving
involution~\cite{tierney1998note}), whereas the IQP refinement is a
many-to-one, cut-decreasing contraction onto local minima, the
opposite of an involution. Nor can the refinement be folded into the
proposal and Metropolized, since the reverse move would require the
argmin at $P'$ to output the strictly worse plan $P$, an event of
probability zero. The distributional signature of this bias is visible
in our own ground-truth data (\S\ref{sec:exp-5x5}), where Chain C's cut
marginal is shifted and sharply concentrated relative to Chain A's,
despite the identical outer kernel.

We therefore present Chain C as what it is, a \emph{matheuristic
ensemble generator} in the lineage of ReCom-based
optimization~\cite{cannon2023voting,brous2025recom}, an iterated
local search whose stochastic component (the tempered outer move)
maintains diversity while the embedded MIQP drives compactness. This
is a legitimate and, as \S\ref{sec:experiments} shows, empirically
powerful tool for producing high-quality compact plans and for
seeding or stress-testing ensembles. What it does not provide is a
defensible null distribution for outlier analysis; for that use case
the field's requirements are exactly the known-target
property~\cite{cannon2022spanning,mccartan2023sequential} that
Chain E (\S\ref{sec:qp-mh}) preserves and Chain C deliberately trades
away for optimization power.

\paragraph{Trade-off: continuous vs.\ integer.}\label{sec:iqp-pareto}
The continuous relaxation is roughly $45\times$ cheaper per step in
wall-clock (a linear solve in the common case, by
Lemma~\ref{lem:row-stochastic}), but on Iowa only the exact integer
solve moves the compactness needle; the contrast is developed in
\S\ref{sec:exp-compactness}.

%════════════════════════════════════════════════════════════════════════════════
\section{Experiments}\label{sec:experiments}
%════════════════════════════════════════════════════════════════════════════════

We evaluate the six chains of Table~\ref{tab:five-chains} on three
datasets, namely the $5 \times 5$ grid (full enumeration, ground-truth
target), the Iowa county dual graph (real geography, 99 counties,
$k = 4$, $\varepsilon = 0.05$), and the North Carolina precinct dual
graph ($2\,692$ precincts, $k = 4$, $\varepsilon = 0.05$). All hyperparameters are
locked in Appendix~\ref{app:hyperparams}; raw cut traces, per-step
balance audits, and configuration stamps are released in the project
repository.

\paragraph{Statistical protocol.} Chains are autocorrelated, so we do
not report i.i.d.\ test statistics on raw traces. For every
comparison we report the mean difference, Cohen's $d$ on the pooled
post-burn-in standard deviation, and a $z$-test whose standard errors
use effective sample sizes (ESS) from Geyer's initial-positive-sequence
estimator of the integrated autocorrelation time. ESS values are
reported alongside and are small, typically tens per 400 retained
Iowa steps (as low as 3 for Chain E). On the $5 \times 5$ grid Chain
C's trace is nearly constant and the estimator is unstable there, so
we draw no single-run inference from it and let the ten-seed
comparisons carry that claim. All ESS values are computed from cut
traces and are appropriate for inference about cut means only;
distributional tests use batch means (\S\ref{sec:exp-2x3}).
Each chain runs on its own seeded RNG stream. The tables below report
a single seed (42) with full diagnostics; every headline comparison
was replicated across ten seeds (42 through 51) under one solver
configuration. On Iowa, Chain C's gain is $-8.0 \pm 0.4$ cuts versus
Chain A ($-20.3\%$), negative on all ten seeds (paired
$p = 6.4 \times 10^{-13}$), with C's mean cut count at
$31.4 \pm 0.3$. Chain B's effect resolves to $-0.9 \pm 0.9$ cuts (population
standard deviation, as in every table)
($-2.3\%$), negative on nine of ten seeds and nominally significant
($p = 0.015$) but an order of magnitude below C's, so single runs
cannot distinguish it from zero. Chain D is indistinguishable from A
($+0.4$ cuts, $p = 0.10$). On the $5 \times 5$ fiber, C's between-seed
spread is essentially nil ($16.03 \pm 0.02$ mean cuts) and its
$-4.0\%$ gain holds on all ten seeds ($p = 4.3 \times 10^{-7}$),
while B is null ($p = 0.59$). Between-seed standard deviations of the
baselines are about $0.4$ cuts on Iowa, so single-seed means are
representative.

\paragraph{Balance audit.}\label{sec:exp-balance} Every experiment
records, at every step, the maximum relative deviation of any district
population from the global ideal. All six chains satisfy the
$\varepsilon = 0.05$ bound at every step of every run reported here
(Iowa maxima 0.0494--0.0500 across the headline and ablation runs). We flag this because it is not automatic.
GerryChain's raw \texttt{recom} proposal validates only one side of
each tree cut, and an earlier version of our own local steps anchored
balance to the merged-pair half rather than the global ideal
(\S\ref{sec:qp-model}); either slip lets district populations drift
to $\sim$4$\times$ the nominal tolerance on Iowa before the audit
catches it. We recommend per-step balance auditing as standard
practice for any custom chain harness that bypasses GerryChain's
validator pipeline.

\subsection{Ground truth on the enumerated fiber: what the family samples}\label{sec:exp-5x5}

The $5 \times 5$ partition fiber (the set of all valid plans; $k = 5$, exact balance, 4\,006
partitions; \S\ref{sec:fiber-corridors}) admits full enumeration, so
every distributional quantity is exactly computable. The chains run
with $\varepsilon = 0.05$, which on 25 unit-population cells admits
only districts of exactly five cells, so the state space the chains
traverse is precisely this enumerated fiber. Uniform-on-fiber is not the target of any chain here; the
total-variation column is a descriptive concentration measure only,
and with about 900 correlated draws over 4\,006 states even an
independent uniform sampler would show TV above $0.78$. The
informative comparisons are the cut means against the exact target
mean and the long runs of \S\ref{sec:exp-longrun}, where distances
are read against proper noise floors.

\begin{table}[H]
\centering
\small
\begin{tabular}{lrrrr}
\toprule
Chain & Mean cuts & Std & Unique / 4006 & TV(emp, uniform) \\
\midrule
V: Vanilla ReCom     & 17.06 & 1.08 & 337 / 4006 (8.4\%) & 0.92 \\
A: Tempered ReCom    & 17.02 & 1.15 & 311 / 4006 (7.8\%) & 0.92 \\
B: QP-diffusion mix  & 16.54 & 0.83 & 153 / 4006 (3.8\%) & 0.96 \\
C: IQP local         & 16.03 & 0.19 & 130 / 4006 (3.2\%) & 0.97 \\
D: Spectral mix      & 16.53 & 0.85 & 208 / 4006 (5.2\%) & 0.95 \\
E: QP-MH exact       & 16.27 & 0.44 & \phantom{00}6 / 4006 (0.1\%) & 1.00 \\
\bottomrule
\end{tabular}
\caption{$5 \times 5$ short-run comparison ($k = 5$, exact balance,
  1\,000 steps, burn-in 100, seed 42). The exact target mean is
  $\E_\pi[|\partial P|] = 18.56$ at $\lambda = 0.10$, and \emph{every}
  chain sits far below it at this horizon, including the exactly
  reversible Chain E, which at 1\,000 steps has visited only 6 states. Exactness of the kernel does not buy mixing speed. Compactness
  gains trade against coverage. Chain C covers 3.2\% of the fiber against
  V's 8.4\%.}\label{tab:5x5-headline}
\end{table}

\textbf{The origin of the gap.} Chains B and C
sacrifice roughly half of V's fiber coverage at this horizon. The
gap between every chain's empirical mean and the target mean 18.56
quantifies the compactness bias \emph{shared} by the whole
ReCom-derived family, and on the enumerated fiber we can identify its
origin. Using the Matrix--Tree theorem we compute, for each of the
4\,006 partitions, the spanning-tree-count weight
$\prod_i \tau(V_i)$, the measure that the RevReCom analysis
identifies as the natural target of tree-based
recombination~\cite{cannon2022spanning,autry2023metropolized} and
whose relation to compactness Procaccia and
Tucker-Foltz~\cite{procaccia2022compact} proved. At the level of
means the match is close. The tree-weighted measure has
$\E[|\partial P|] = 16.88$ against Vanilla ReCom's ten-seed mean of
$16.95 \pm 0.11$, and the tree-weighted-and-tempered measure
$\prod_i \tau(V_i)\, e^{-\lambda|\partial P|}$ has
$\E[|\partial P|] = 16.78$ against Chain A's ten-seed mean of
$16.71 \pm 0.16$; both gaps are 0.07 cuts, within one between-seed
standard deviation, and the two quasi-targets sit 1.7 and 1.8 cuts
below the nominal tempered mean. A mean is a one-dimensional summary,
though, and many distributions share one. The long runs below compare
full distributions, which is where the claim is actually tested. Our
refinements deepen the pre-existing bias only modestly (Chain C lands
2.5 cuts below the nominal target). Chain C's relative gain is
far larger on Iowa ($-20.3\%$ across ten seeds,
\S\ref{sec:exp-compactness}, against $-4.0\%$ here), reflecting the
room that heterogeneous county populations leave for optimization;
the uniform grid leaves little, so we read the $5 \times 5$ results
as calibration, not as a demonstration of power.

\paragraph{Long-run validation at the level of distributions.}\label{sec:exp-longrun}
We run Chain A for $2 \times 10^5$ steps and Chain E for $10^5$ steps
on the enumerated fiber (seed 42, 10\% burn-in) and compare each
chain's empirical distribution against three exactly computable
reference measures, the nominal target
$\pi \propto e^{-\lambda|\partial P|}$, the tree-weighted measure
$\prod_i \tau(V_i)$, and the tree-weighted-and-tempered measure. Two
design choices make the comparison informative. First, a raw
total-variation distance is uninterpretable on 4\,006 states, since
even an independent sampler shows large TV until the draw count
greatly exceeds the state count. So every distance is reported
against a parametric-bootstrap noise floor, the same distance
computed for independent draws from the reference measure at the
chain's effective sample size. A converged chain sits at its floor; a
chain converging elsewhere sits above it by many floor standard
deviations. Second, the headline statistic is the cut-count marginal
rather than the full state distribution. The fiber's cut counts take
five values, so with $1.8 \times 10^5$ retained draws the marginal is
estimated to three decimals and its floor is small, which removes the
coverage problem entirely.

\begin{table}[H]
\centering
\small
\begin{tabular}{llrrr}
\toprule
Chain & Reference measure & TV & Floor at ESS & $z$ \\
\midrule
A ($2\times10^5$ steps, ESS 12\,312) & nominal $\pi$ & 0.629 & 0.007 & $+220$ \\
 & tree-weighted & 0.065 & 0.006 & $+22$ \\
 & tree-weighted, tempered & 0.022 & 0.006 & $+6.7$ \\
\midrule
E ($10^5$ steps, ESS 28) & nominal $\pi$ & 0.079 & 0.141 & $-1.2$ \\
 & tree-weighted & 0.493 & 0.125 & $+6.9$ \\
 & tree-weighted, tempered & 0.532 & 0.125 & $+7.5$ \\
\bottomrule
\end{tabular}
\caption{Total variation between each chain's empirical cut-count
  marginal and three reference measures on the enumerated
  $5 \times 5$ fiber, with the parametric-bootstrap noise floor at the
  chain's effective sample size and the excess in floor standard
  deviations. Chain A is far from its nominal target and close to,
  though distinguishable from, the tree-weighted tempered measure.
  Chain E is indistinguishable from the nominal target it is proved
  reversible for and far from the tree measures.}\label{tab:calibration-tv}
\end{table}

\begin{figure}[H]
  \centering
  \includegraphics[width=\textwidth]{fig/calibration_5x5.pdf}
  \caption{What the chains sample. (a)~Chain A's cut-count marginal
    after $2 \times 10^5$ steps against the three reference measures.
    (b)~Total variation of the running marginal to each reference as a
    function of run length, with the noise floor at the effective
    sample size shaded; the distance to the nominal target is a
    plateau across five decades while the distance to the
    tree-weighted tempered measure descends toward the floor.
    (c)~Chain E's marginal after $10^5$ steps against its own
    target.}
  \label{fig:calibration}
\end{figure}

\emph{Chain A does not sample its nominal target, and the gap is a
plateau, not a drift.} At $2 \times 10^5$ steps the tempered chain
has visited $3\,524/4\,006$ states, its mean is $16.73$, and its cut
distribution sits 220 floor standard deviations from the nominal
target, at TV $0.63$. Figure~\ref{fig:calibration}(b) shows that
distance flat across five decades of run length while the distance
to the tree-weighted tempered measure falls toward the floor. This is
convergence to a different distribution, not slow mixing, and it is
what the symmetric-$q$ filter of \S\ref{sec:qp-mh} predicts. The
match to the tree-weighted tempered measure is close but not exact.
At TV $0.022$ the chain sits $6.7$ floor standard deviations above
it, so the tree measure is a very good description of what tempered
ReCom samples and not a perfect one, which is expected, since Vanilla
ReCom's stationary law is not exactly the tree measure and RevReCom
exists to make it so. On the full 4\,006-state distribution the noise
floor at effective sample size is $0.15$ to $0.23$, so that comparison
can separate the nominal target ($z = +164$) but not the two tree
measures from each other, which is why the marginal carries the
claim.

\emph{Chain E samples what it states.} At $10^5$ steps its cut
marginal is $1.2$ floor standard deviations from the nominal target
and $7$ or more from either tree measure. The two chains separate
exactly as theory predicts. Convergence of the full distribution is
far slower for this sticky local chain, as \S\ref{sec:mobility}
explains. It has visited $403/4\,006$ states, and its full-state TV
of $0.90$ sits at a floor of $0.99$ at effective sample size 28, so
that test has no power here. The fiber on which full convergence can
be checked is the next subsection.

\subsection{Exact convergence on the 18-state fiber}\label{sec:exp-2x3}

Detailed balance gives stationarity, not convergence, and on 4\,006
states the exact chain cannot be run to convergence in reasonable
time. On the $2 \times 3$ grid with $k = 2$ and tolerance
$\varepsilon = 1/3$, the fiber of labeled valid plans has 18 states,
the same state space on which Remark~\ref{rem:qpmh-scope} verifies
the kernel to machine precision, and a long run can cover it. We run
Chain E for $10^6$ steps (seed 42, $10\%$ burn-in, locked parameters
$\alpha = 1$, $\beta = 10$, $\lambda = 0.10$). It visits all 18
states, its acceptance rate is $0.785$, and its post-burn-in mean cut
count is $2.75992$ against the exact $2.76002$.

The distributional test uses batch means, because thinning a chain by
an effective sample size computed from a summary statistic is not
valid here. The cut count takes two values on this fiber, so its
integrated autocorrelation time ($17$ steps) understates that of the
18-state occupancy process ($25$ steps), and a chi-square test on a
chain thinned by the former rejects spuriously. We therefore split
the trace into 200 blocks of $4\,500$ steps, over 180 occupancy
correlation times each, estimate the covariance of the occupancy
vector from the between-block scatter, and test with Hotelling's
$T^2$ on the 17 free coordinates. The result is $F(17, 183) = 0.87$,
$p = 0.61$. The occupancy distribution is not distinguishable from
the exact target. Total variation to the target is $0.0110$, against
a noise floor of about $0.009$ at the occupancy effective sample
size, and it fell from $0.0250$ at $2 \times 10^5$ steps, a factor of
$2.27$ for five times the steps, which is the $\sqrt{n}$ rate of a
converged chain. At $2 \times 10^5$ steps the same test gave
$p = 0.07$ with three states in one cut-level set persistently
under-visited; at $10^6$ steps those deviations dissolved into the
one or two $|z| \approx 2$ values that 18 states produce by chance.
Together with the kernel check, this is end-to-end validation of the
construction of \S\ref{sec:qp-mh}. The implemented kernel satisfies
detailed balance exactly, and the chain it generates reaches its
target.

\subsection{Compactness on Iowa: refinement gains and optimizer baselines}\label{sec:exp-compactness}

\paragraph{Iowa (6 chains).}\label{sec:exp-iowa}

\begin{table}[H]
\centering
\small
\begin{tabular}{lrrllr}
\toprule
Chain & Mean cuts & Std & $\Delta$ vs A (Cohen's $d$) & ESS-$z$ $p$ & Wall-clock \\
\midrule
V: Vanilla ReCom         & 41.39 & 4.78 & ---                & ---     & 1.2\,s \\
A: Tempered ReCom        & 39.76 & 4.32 & ---                & ---     & 1.1\,s \\
B: QP-diffusion mix      & 38.15 & 4.28 & $-4.0\%$ ($-0.37$) & $0.24$  & 1.3\,s \\
\textbf{C: IQP local ($n_{\mathrm{iqp}} = 10$)}
                         & \textbf{31.78} & \textbf{1.25}
                         & \textbf{$-20.1\%$ ($-2.51$)} & $< 10^{-11}$ & 57.2\,s \\
D: Spectral mix          & 39.51 & 4.70 & $-0.6\%$ ($-0.06$) & $0.85$  & 3.9\,s \\
E: QP-MH exact           & 55.21$^{\dagger}$ & 7.88 & n/a (not converged) & n/a & 3.6\,s \\
\bottomrule
\end{tabular}
\caption{Iowa cut counts ($k = 4$, $\varepsilon = 0.05$, 500 steps,
  burn-in 100, seed 42, per-step balance audit passed by all chains).
  Chain C delivers a 20.1\% reduction in mean cuts versus Tempered
  ReCom on this seed ($-20.3\%$ across ten seeds) with the tightest
  distribution; Chains B and D are
  statistically indistinguishable from A once balance is enforced
  correctly and autocorrelation is accounted for. $^{\dagger}$Chain E
  is not a competitor on compactness. It is the exact sampler for the
  tempered target, and its trajectory is still rising at 500 steps
  (ESS $\approx 3$), so we report its mean as a \emph{non-equilibrium
  lower bound} on the target's cut level, with no test statistic.
  Even as a lower bound it shows the target
  $\pi \propto \exp(-0.1\,|\partial P|)$ is far \emph{less} compact
  than the output of any ReCom-family chain.}\label{tab:iowa-headline}
\end{table}

Three readings of Table~\ref{tab:iowa-headline}. First, the
integer refinement is the only local kernel whose Iowa gain survives
rigorous accounting. A change of $-20.1\%$ on this seed ($-20.3\%$
across ten seeds) is a large difference within the chain family (the
reported $d = -2.51$ scales it by the chain's equilibrium spread,
which is not sampling noise, and should be read descriptively),
though
\S\ref{sec:exp-optimizer-baselines} shows classical annealing reaches
the same best plans far more cheaply, and it is robust to the
balance-constraint correction that erased Chain B's
previously-measured advantage (with the global window enforced, B's
point estimate swings between $-0.3\%$ and $-4.0\%$ across equivalent
reruns and never approaches significance at these effective sample
sizes; \S\ref{sec:disc-limits}).
Second, Chain D's parity with A and V supports our claim that the
spectral cut alone, without the fidelity term, does not improve
ReCom-family compactness on heterogeneous county graphs. Third, Chain
E's trajectory shows the Iowa version of the calibration result of
\S\ref{sec:exp-5x5}. Taking E's non-equilibrium mean as a lower
bound, every ReCom-based chain here sits at least $\sim$14--23 cuts
\emph{below} the tempered target's own mean, because the
spanning-tree proposal is itself strongly compactness-biased.
``Tempered ReCom targets $\pi$'' is, on this evidence, a statement
about the filter, not the samples.

\begin{figure}[H]
  \centering
  \includegraphics[width=0.92\textwidth]{fig/target_matched_iowa.png}
  \caption{Cut-edge trace and post-burn-in histogram for the six
    chains on Iowa. Chain C concentrates tightly at low cut counts;
    Chain E drifts upward toward the tempered target's own cut
    level.}
  \label{fig:iowa-trace}
\end{figure}

The wall-clock column completes the practitioner picture. Chain C's
compactness costs $\sim 47\times$ the per-step time of Vanilla ReCom at
$n_{\mathrm{iqp}} = 10$ (about 0.11\,s per step on Iowa), which is
acceptable for ensembles of thousands of plans but material at
precinct scale. Chain C's tight standard deviation (1.25 vs.\ 4.32
for A) is the signature of an optimizer concentrating on IQP local
minima rather than the spread of a sampler at equilibrium, consistent
with its matheuristic reading (\S\ref{sec:chain-c-honest}).

\paragraph{Optimizer baselines.}\label{sec:exp-optimizer-baselines}
Because Chain C is benchmarked as an optimizer, we compare it against
the classical optimizers a referee should demand. Simulated annealing
over ReCom proposals (the same Metropolis filter as Chain A with a
geometric cooling schedule, $\lambda$ from $0.05$ to $5$) and pure
hill climbing (accept only non-worsening moves) both cost the same as
Vanilla ReCom, about one second for 500 steps on Iowa. We also run
the same schedule stretched to $5\,000$ steps and a 25-restart
variant.

\begin{table}[H]
\centering
\small
\begin{tabular}{lrrrrr}
\toprule
Method & Best cuts & Ensemble mean & Distinct plans & Pairwise dist. & Wall \\
\midrule
SA (500 steps)           & 31 & 31    & 1 (frozen)  & ---  & 1.0\,s \\
Hill climbing (500)      & 31 & 31    & 1 (frozen)  & ---  & 0.9\,s \\
SA (5\,000 steps)        & 29 & 31    & 1 (frozen)  & ---  & 9.7\,s \\
SA, 25 restarts          & 29 & 31.00 & 23          & 25.1 & 26\,s \\
Vanilla ReCom + IQP      & 31 & 31.68 & moving      & ---  & 53\,s \\
Chain C                  & 29 & 32.34 & 135         & 23.5 & 57\,s \\
\bottomrule
\end{tabular}
\caption{Optimizer baselines on Iowa (seed 42,
  Gurobi 13). ``Ensemble mean'' is the mean cut count of the plans
  each method ends up producing (the frozen plateau for single-run
  SA and hill climbing, the set of 23 distinct final plans for
  restarted SA, the distinct post-burn-in plans for the chain
  methods). ``Distinct plans'' counts distinct plans produced (over
  400 post-burn steps for Chain C). ``Pairwise dist.''\ is the mean
  number of counties reassigned between two distinct plans under the
  best district relabeling. The single-run SA and hill-climbing rows
  were replicated across ten seeds. SA's best ranges over 30--32 and
  never reaches 29 in 500 steps (mean frozen plateau 31.8), and hill
  climbing lands at 31 on every
  seed.}\label{tab:optimizer-baselines}
\end{table}

We report this comparison as a main finding rather than bury it, and
it cuts both ways. The best plan quality found (29 cut edges) is
reached by three routes, namely the $5\,000$-step anneal (9.7\,s),
restarted annealing (26\,s), and Chain C (57\,s). Single 500-step
anneals stop at 30--32 across ten seeds and hill climbing at 31 on
every seed, so the compact tail takes patience or restarts. For a
\emph{single} compact plan, the longer anneal is the clear winner at
about a sixth of Chain C's cost, so nobody should reach for an MIQP
to find one compact map. As \emph{ensemble} generators the two are
closer. At under half of Chain C's wall-clock, restarted annealing yields a
slightly better mean over its distinct plans (31.00 vs.\ 32.34) and
slightly higher diversity (25.1 vs.\ 23.5 counties reassigned between
plan pairs), while Chain C yields far more distinct near-optimal
plans per minute (about 141 vs.\ 53). Neither dominates. A
$2 \times 2$
decomposition sharpens where the compactness comes from. Adding IQP
refinement to Vanilla ReCom is worth about $10$ cuts on the mean
(41.39 to 31.68) and adding it to Tempered ReCom about $8$ (39.76 to
31.78), while the tempering filter alone is worth about $1.6$ cuts
and moves the post-refinement mean by under $0.2$. The optimizer does
the work; the outer filter barely moves the mean, though on this seed
only the tempered variant reached a 29-cut plan.

Two conclusions follow. On pure compactness optimization the classical
baseline matches the MIQP machinery at a fraction of the cost, and we
believe papers in the optimization-inside-ReCom line should report
it as standard practice. What the energy retains is anchoring, small
changes near a given plan, tested head to head in
\S\ref{sec:exp-leastchange}, together with kernel flexibility
(\S\ref{sec:exp-multikernel}) and, in Chain E, a provable stationary
distribution that no annealing schedule provides.

\subsection{North Carolina precinct validation}\label{sec:exp-nc}
We repeat the six-chain comparison on the North Carolina precinct dual
graph (\texttt{NC\_VTD}, $2\,692$ precincts, $k = 4$,
$\varepsilon = 0.05$, seed 42, 500 steps), roughly twenty-seven times
the vertex count of the Iowa instance. This is a single-seed scale
demonstration, so we report effect sizes without ESS-corrected test
statistics and read the directions, not the exact magnitudes.

\begin{table}[H]
\centering
\small
\begin{tabular}{lrrlr}
\toprule
Chain & Mean cuts & Std & $\Delta$ vs A (Cohen's $d$) & Wall-clock \\
\midrule
V: Vanilla ReCom    & 302.5 & 37.9 & ---                & 31\,s \\
A: Tempered ReCom   & 226.4 & 23.8 & ---                & 30\,s \\
B: QP-diffusion mix & 252.6 & 50.9 & $+11.5\%$ ($+0.66$) & 89\,s \\
\textbf{C: IQP local ($n_{\mathrm{iqp}} = 10$)}
                    & \textbf{145.7} & \textbf{3.9}
                    & \textbf{$-35.7\%$ ($-4.73$)} & 1950\,s \\
D: Spectral mix     & 253.7 & 49.8 & $+12.0\%$ ($+0.70$) & 131\,s \\
E: QP-MH exact      & 418.6$^{\dagger}$ & 25.1 & n/a (not converged) & 91\,s \\
\bottomrule
\end{tabular}
\caption{North Carolina precincts ($k = 4$, $\varepsilon = 0.05$,
  500 steps, burn-in 100, seed 42, per-step balance audit passed).
  Chain C reduces mean cut edges by $35.7\%$ versus Tempered ReCom,
  a larger relative gain than on Iowa, at a per-step cost of about
  $3.9$\,s. $^{\dagger}$Chain E is still climbing at 500 steps and its
  mean is a non-equilibrium lower bound, as on Iowa.}\label{tab:nc-headline}
\end{table}

Two things carry over from Iowa and one is new. The integer refinement
again dominates its chain family, and by a wider margin
($-35.7\%$, $d = -4.73$), consistent with our reading that
heterogeneous real-world graphs leave more room for local
optimization than the uniform grid. Chain C's distribution is again
sharply concentrated (std $3.9$ against $23.8$ for A). The new
observation is that on this larger graph the two mixture chains, B and
D, land \emph{above} Tempered ReCom rather than beside it ($+11.5\%$
and $+12.0\%$), because a single continuous QP solve or spectral cut
on a precinct-scale merged region is a coarse move that the
unconditional accept step does not filter, so at scale the cheap
regimes hurt. Only the exact integer solve (Chain C) and the exact
sampler (Chain E, by construction) behave well at precinct scale. The
cost is real. Chain C's $1950$\,s for 500 steps, about $3.9$\,s per
step, is acceptable for a targeted refinement pass but not for a long
ensemble run, which is the same speed-versus-power trade the Iowa
timings show, magnified by scale.

\subsection{Multi-kernel showcase on Iowa}\label{sec:exp-multikernel}

We run Chain B and Chain C on Iowa with three kernels and visualize
the resulting partitions. The three kernels are uniform
($k_{\mathrm{unif}}$), perimeter ($k_{\mathrm{perim}}$,
$w_{ij}$ = shared-boundary length from \texttt{shared\_perim}), and
density ($k_{\mathrm{dens}}$, $w_{ij} = |\partial v_i \cap \partial v_j| \cdot \exp(-(\rho_i - \rho_j)^2 / \sigma^2)$
with $\rho_v = \mathrm{TOTPOP}/\mathrm{area}$ normalized to unit
mean and $\sigma = 1$).

\begin{figure}[H]
  \centering
  \includegraphics[width=\textwidth]{fig/multikernel_iowa.png}
  \caption{Final Chain C plans on Iowa under the three kernels
    (seed 42, 500 steps, identical constraints; only the Laplacian
    edge weights differ). Panel annotations give the unweighted cut
    count and the geographic boundary length (sum of shared-perimeter
    weights over cut edges). The plans differ in 19 to 25 of the 99
    counties under optimal district relabeling. The kernel directly
    encodes geographic signal at the proposal level.}
  \label{fig:multikernel-iowa}
\end{figure}

\begin{table}[H]
\centering
\small
\begin{tabular}{lrrrr}
\toprule
Kernel & B mean cuts & C mean cuts & C boundary (mean $\pm$ sd)
  & C $|\Delta\rho|$ ratio (mean $\pm$ sd) \\
\midrule
$k_{\mathrm{unif}}$  & 38.15 & 31.78 & $10.38 \pm 0.67$ & $1.78 \pm 0.36$ \\
$k_{\mathrm{perim}}$ & 39.00 & 33.39 & $\mathbf{9.89 \pm 0.42}$ & $1.82 \pm 0.39$ \\
$k_{\mathrm{dens}}$  & 38.09 & 36.84 & $11.31 \pm 1.16$ & $\mathbf{2.07 \pm 0.37}$ \\
\bottomrule
\end{tabular}
\caption{Three kernels on Iowa ($k = 4$, $\varepsilon = 0.05$,
  500 steps, burn-in 100, seed 42, corrected balance constraint,
  Gurobi 13). Mean cuts are post-burn-in averages of the unweighted
  cut count. The last two columns are averaged over every post-burn-in
  Chain C plan (not one snapshot), with the standard deviation across
  plans. ``C boundary'' is the geographic boundary length (sum of
  \texttt{shared\_perim} over cut edges, in the dual graph's units).
  ``$|\Delta\rho|$ ratio'' is the mean population-density gap across
  cut edges divided by the graph-wide mean gap, so $1.00$ means
  density-blind boundaries and larger values mean boundaries
  concentrated on urban--rural transitions. Each kernel is best on its
  own metric (bold): the perimeter kernel gives the shortest
  geographic boundary, the density kernel the highest density-gap
  ratio. Under non-uniform kernels the unweighted cut count is not the
  objective, so the cut columns must be read together with the
  kernel-specific columns.}\label{tab:multikernel-iowa}
\end{table}

\textbf{Positioning of the kernel showcase.} SpecReCom uses randomized
edge weights only to break spectral ties; region-aware ReCom and
county-weighted trees accept indicator surcharges; none accepts an
arbitrary user-supplied kernel with a locality anchor. The claim this
experiment supports is deliberately scoped to kernels as first-class
proposal inputs, with characterized behavior, not priority over
all weighted proposals (\S\ref{sec:related-weighted}).

\textbf{What each kernel does to the plan.} The three final Chain C
plans in Figure~\ref{fig:multikernel-iowa} are genuinely different
districtings, not perturbations of one another. Under the best
district relabeling they disagree on 19 of 99 counties (uniform vs.\
perimeter), 25 (uniform vs.\ density), and 24 (perimeter vs.\
density). Averaged over the whole post-burn-in trajectory, each
kernel is best on the quantity it encodes. The perimeter kernel
prices each edge by its shared boundary length and gives the shortest
mean geographic boundary ($9.89$ vs.\ $10.38$ for uniform, a $4.7\%$
reduction). The density kernel prices edges by shared boundary length
damped across density gaps, and it does the opposite on raw boundary,
producing the \emph{longest} average boundary ($11.31$), because it
will accept a longer border in exchange for aligning cuts with
population-density transitions. Its mean density gap across cut edges
is $2.07\times$ the graph-wide mean, against $1.78\times$ for the
uniform kernel. In map terms, the density plan's four largest-gap cut
edges are precisely the borders of Polk County (Des Moines), followed
by those of Linn (Cedar Rapids) and Johnson (Iowa City), the state's
three main urban cores, rather than edges through uniformly rural
blocks.

We stress the trajectory averages over single-plan snapshots on
purpose. A single final plan is a noisy draw, and its boundary and
gap numbers can point the wrong way. Figure~\ref{fig:multikernel-iowa}
itself illustrates this. The final density-kernel plan happens to land
a shorter boundary than the final perimeter plan, opposite to the
trajectory means. Only the post-burn-in means, with the
spreads in Table~\ref{tab:multikernel-iowa}, separate the kernel's
systematic effect from run-to-run noise.

\textbf{Reading the cut-count column.} Under non-uniform
kernels the unweighted cut count worsens, from 31.78 (uniform) to
33.39 (perimeter) and 36.84 (density) for Chain C, and it should, because
raw cut count is no longer the objective. The kernel-specific metrics
improve instead, which is the point of the showcase. Chain B is
nearly kernel-insensitive (38.1--39.0 mean cuts across all three),
consistent with its Iowa null result in
Table~\ref{tab:iowa-headline}. The continuous relaxation plus
rounding does not transmit the kernel's signal into the final integer
plan the way the exact integer solve does.

\subsection{Least-change refinement from the enacted plan}\label{sec:exp-leastchange}

Remedial redistricting often asks for the smallest change that
improves a specific plan, for example when a commission must repair a
defect without redrawing the whole map. This is the setting the
fidelity term was built for, and we quantify it directly. The
reference plan is Iowa's enacted 2011 congressional map (47 cut
edges; district populations within $0.01\%$ of the ideal, so it is a
valid state of every chain in this paper). A method may reassign at
most $D$ of the 99 counties, counted under the best district
relabeling, and we record the best cut count it achieves within that
budget, over ten seeds.

Three methods start at the reference. The \emph{anchored IQP
refinement} is the paper's inner move alone ($\alpha = 10$, uniform
kernel, global balance window, no outer ReCom move), run for 50
passes. Because the fidelity weight $\beta$ is exactly the
least-change dial, the method's frontier is the envelope of (cuts,
distance) trajectories over $\beta \in \{0, 0.3, 1, 3, 10, 30\}$,
about 3\,s of total solve time per seed. \emph{Distance-capped
annealing over ReCom proposals} uses the geometric schedule of
\S\ref{sec:exp-optimizer-baselines} and rejects any proposal beyond
the budget ($2\,000$ proposals per budget, about 4\,s).
\emph{Distance-capped annealing over single-county boundary flips},
with contiguity and the global balance window enforced per flip, is
the fine-grained classical alternative ($2\,000$ proposals, under
1\,s).

\begin{table}[H]
\centering
\small
\begin{tabular}{rrrr}
\toprule
Budget $D$ & IQP anchored ($\beta$ envelope) & SA-ReCom capped & SA-flip capped \\
\midrule
3  & $42.40 \pm 1.36$ & $41.40 \pm 0.80$ & $\mathbf{40.60 \pm 0.66}$ \\
5  & $39.60 \pm 0.92$ & $38.80 \pm 0.75$ & $\mathbf{38.20 \pm 0.40}$ \\
10 & $\mathbf{32.00 \pm 0.00}$ & $32.10 \pm 0.30$ & $33.10 \pm 0.94$ \\
15 & $32.00 \pm 0.00$ & $\mathbf{31.70 \pm 0.46}$ & $32.10 \pm 0.30$ \\
20 & $32.00 \pm 0.00$ & $\mathbf{31.20 \pm 0.40}$ & $31.60 \pm 0.49$ \\
none & $31.80 \pm 0.40$ & $\mathbf{30.60 \pm 0.66}$ & $32.30 \pm 0.46$ \\
\bottomrule
\end{tabular}
\caption{Least-change refinement of the enacted Iowa plan (47 cut
  edges). Entries are the best cut count reachable while reassigning
  at most $D$ counties (best district relabeling), mean $\pm$ sd over
  ten seeds; bold marks the best method per budget. The last row
  removes the cap ($D = 99$).}\label{tab:least-change}
\end{table}

Move granularity explains the pattern. From the reference, an
accepted ReCom move reassigns a median of 26 counties (mean 25.0,
minimum 1 over 197 accepted moves), so the capped ReCom annealer
survives on the rare tree draws that happen to re-split almost
identically. An accepted IQP pass at $\beta = 1$ moves 2.5 counties
on average (median 2, maximum 5).

No method dominates the frontier, and we report all three regimes.
At tight budgets ($D \le 5$) capped flip annealing wins; one-county
moves explore the small-perturbation space that both the tree
re-split and the pair-optimal IQP step overshoot. At $D = 10$ the
anchored refinement ties the capped ReCom annealer ($32.00$ against
$32.10 \pm 0.30$, inside one standard deviation) and, unlike the
annealers, it is deterministic. From the enacted plan the $\beta = 1$ refinement
removes 15 of 47 cut edges while moving exactly 10 counties, reaching
the same plan on all ten seeds in about half a second of solve time,
with no cooling schedule or restart budget to tune. At loose budgets
($D \ge 15$) the task stops being least-change and capped ReCom
annealing pulls ahead, consistent with
\S\ref{sec:exp-optimizer-baselines}. The practical summary is that
for repairs moving roughly ten percent of units the anchored IQP
matches the best annealer and is the only reproducible method,
smaller repairs favor flip annealing, and larger redraws favor global
annealing. An exact benchmark for this setting exists in the
transfer-distance midpoint programs of Dobbs et
al.~\cite{dobbs2024facilitating}, and comparing against it is the
natural next step for this experiment. The IQP
frontier here came from sweeping $\beta$, which confirms the fidelity
term as a usable least-change dial rather than a fixed setting.

%════════════════════════════════════════════════════════════════════════════════
\section{Discussion}\label{sec:discussion}
%════════════════════════════════════════════════════════════════════════════════

\subsection{What each member of the family is for}\label{sec:disc-roles}

The three new chains occupy different points on the
correctness--power spectrum, and no single point dominates.
\textbf{Chain C} delivers a large, replicated compactness gain over
its own chain family, but that is an optimizer's power, and annealed
ReCom reaches the same best plans within seconds
(\S\ref{sec:exp-optimizer-baselines}). Its niche is anchored,
kernel-aware refinement of a given plan at moderate change budgets,
where it matches the best annealer and is the only reproducible
method (\S\ref{sec:exp-leastchange}); for compactness search use
annealing with restarts, and use neither to build a null distribution
for outlier analysis. \textbf{Chain E} is the opposite pole,
provably exact for its target and, by the same token, unsparing about
what that target looks like. Its Iowa trajectory climbing past 55
cuts, with every other chain at least $14$ to $23$ cuts below, shows
on real geography what the enumerated fiber shows exactly, that the
family's outputs sit far from the tempered distribution its filters
nominally target. \textbf{Chain B} occupies the middle and is a
negative result on Iowa, kept because the QP/IQP contrast shows that
integrality rather than spectral smoothing is doing the work, and
because its per-step cost is $45\times$ lower than C's.

Why does the fidelity term matter? The evidence is the Iowa
compactness contrast, not the jump-rate table. Chain D, our
balance-constrained variant of the SpecReCom idea with no $\beta$
term, is statistically indistinguishable from Vanilla and Tempered
ReCom on Iowa (Table~\ref{tab:iowa-headline}). The spectral cut alone
re-splits each merged pair from scratch and discards the structure the
chain has already found, whereas the anchored integer refinement
improves on it decisively. The community step-trace diagnostic
(Appendix~\ref{sec:fiber-corridors}) makes a complementary, weaker point
about locality in general. Pure tree-move chains cross community
boundaries at 13--15\% per step, near the community-blind baseline,
while every mixture with an optimization-guided local move (anchored
or not) drops the rate to 5--11\%. The hybrid design, ReCom for
across-community mobility and anchored refinement for
within-community descent, follows from the two observations
together.

\subsection{On the macro-structure findings}\label{sec:disc-throats}

The corrected picture of Appendix~\ref{sec:fiber-corridors} has two
consequences. First, it relocates the blame for
slow local mixing. If plan space were islands joined by narrow
corridors, slow traversal would be a topological inevitability and no
local proposal could fix it; since the fibers are in fact well
connected, slow traversal is
a property of the dynamics, and better local proposals are not
hopeless. Chain E crossing a community boundary 3 times in 1\,000
steps despite 18\% of edges being available is the extreme data point. Its
bottleneck is acceptance of large moves, not the graph. Second, the
finding sharpens what a mixing obstruction for these state spaces
would have to look like. Worst-case hardness
results~\cite{akitaya2022reconfiguration,charikar2023complexity}
permit bottlenecked instances, Gold~\cite{gold2026adversarial}
exhibits planar-grid configurations that trap recombination for
exponentially many expected steps, and our grids may simply be too
small or too symmetric to exhibit such structure; enumerating a fiber with a genuine
conductance bottleneck (or proving none exists for grid families)
would be a valuable follow-up, and the statistics of
Table~\ref{tab:communities} give the template for checking.

\subsection{Limitations}\label{sec:disc-limits}

\textbf{(L1) Chain B adds nothing significant on Iowa}
(\S\ref{sec:exp-compactness}); the relaxation appears too loose on
graphs with highly variable vertex populations.

\textbf{(L2) No theoretical mixing-time bound.} Like every existing
redistricting chain, we lack a polynomial mixing-time guarantee, and
Chain E demonstrates that even an exactly reversible member of the
family can be an extremely slow mixer. We rely on empirical
diagnostics, namely ESS, the community step-trace diagnostic, and enumerated
ground truth.

\textbf{(L3) Chain A's symmetric-$q$ approximation is not a small
bias.} On the enumerated fiber the tempered chain's empirical mean
($16.73$) sits $1.8$ cuts below the exact target mean ($18.56$), its
cut distribution is 220 noise-floor standard deviations from that
target (\S\ref{sec:exp-longrun}), and
Chain E's Iowa trajectory suggests the gap is larger, not smaller, on
real graphs. Correcting it requires spanning-tree-count acceptance
ratios~\cite{autry2023metropolized} or a RevReCom outer
kernel~\cite{cannon2022spanning}.

\textbf{(L4) Chain E mixes too slowly to use alone.} Exactness comes
from single-shot rounding, whose reach Corollary~\ref{cor:small-tau}
bounds and whose feasibility \S\ref{sec:mobility} measures; the chain
visited 6 of 4\,006 fiber states in 1\,000 steps at the locked
timestep, which \S\ref{sec:mobility} shows sits well below the
mobility optimum. It validates the corrected
proposal and calibrates other chains' bias; it is not a production
sampler.

\textbf{(L5) Integer QP is computationally expensive.} Chain C costs
$\sim 47\times$ Vanilla ReCom per step at $n_{\mathrm{iqp}} = 10$ on
Iowa. For precinct-scale graphs, time limits with warm starts (and
possibly coarsened merged-pair subproblems) will be necessary; the
NC experiment (\S\ref{sec:exp-nc}) is the pending scalability test.

\textbf{(L6) Detail tables are single-seed}, and should be read with
the ten-seed replication of \S\ref{sec:experiments}.

\subsection{Future directions}\label{sec:disc-future}

\begin{enumerate}
  \item \textbf{A reversible hybrid.} Mix Chain E's exact local
    kernel with a global move that is Metropolized toward the
    \emph{same} target $\pi$ (e.g., a RevReCom-style tree move
    accepted with its computable spanning-tree-count Hastings
    ratio~\cite{cannon2022spanning,autry2023metropolized}), giving a
    chain that is mobile across communities and exactly stationary for a
    stated target. (A naive mixture of kernels reversible for
    \emph{different} targets preserves neither.)
    Corollary~\ref{cor:small-tau} makes this necessary rather than
    optional, since no timestep lifts the local move's reach above
    half a merged pair. The natural target for the hybrid is
    kernel-weighted, $\pi_w(P) \propto \exp(-\lambda\, w(\partial P))$,
    which costs one line in the acceptance ratio, leaves
    Proposition~\ref{prop:qpmh-db} unchanged, and makes the Laplacian
    in the proposal and the weights in the target the same object, so
    that the diffusion move is matched to the target's geometry by
    construction. Metropolized forest
    recombination~\cite{autry2023metropolized} already targets a
    stated law with a tree move, so the question the hybrid must
    answer is what the local kernel adds, and we state the test in
    advance: effective samples per second on $\pi_w$ against the tree
    move alone, at matched wall-clock, on the enumerated
    $5 \times 5$ fiber where convergence is checkable exactly and
    then on Iowa.
  \item Replica exchange or parallel tempering between Chain C and
    sampling chains, in the spirit of multiscale parallel
    tempering~\cite{chuang2025multiscale}, to combine compactness with
    coverage.
  \item Integration with the cycle walk~\cite{deford2025cycle} or the
    balanced up-down walk~\cite{akitaya2026bud} as outer kernels.
  \item Benchmarking Chain C against exact cut-minimization optima on
    the public instances of Validi and
    Buchanan~\cite{validi2022political}, positioning it properly in
    the matheuristic literature~\cite{brous2025recom}.
  \item Multi-scale diffusion kernels (heat kernel, commute-time
    kernel), and a kernel-faithful Cheeger sweep that scores
    candidates by weighted cut (\S\ref{sec:rounding}).
\end{enumerate}

%════════════════════════════════════════════════════════════════════════════════
\section{Conclusion}\label{sec:conclusion}
%════════════════════════════════════════════════════════════════════════════════

This paper introduced a diffusion proposal for redistricting Markov
chains, one Laplacian energy
$\alpha\,x^\top L x + \beta\,\|x - x_0\|^2$ under global population
balance, and showed that the diffusion view is more than an
interpretation. The proposal is a single linear solve whenever
balance is slack, its expected move size is an explicit filter norm
with a small-timestep law and a hard ceiling, and its mobility is
governed by feasibility rather than by the Metropolis ratio, with an
interior optimum in the timestep that the locked setting missed by a
factor of about four. Corrected by an exact Metropolis--Hastings
step, the proposal gives Chain E, to our knowledge the first
optimization-based proposal in redistricting MCMC with a
detailed-balance argument, and on an 18-state fiber a million-step
run's occupancy distribution is not distinguishable from the exact
target. On the enumerated $5 \times 5$ fiber we showed at the level
of distributions that tempered ReCom does not sample its nominal
target and sits close to the spanning-tree-weighted measure, while
Chain E samples what it states. The same energy as an integer program
(Chain C) is a matheuristic that delivers a $20.3\%$ cut reduction
over Tempered ReCom on Iowa across ten seeds, though annealing
matches its best plans at far lower cost, so its role is anchored,
kernel-aware refinement near a given plan; the continuous relaxation
with sweep rounding (Chain B) adds nothing significant on
heterogeneous graphs. Per-step balance and ESS audits discipline the
empirical claims, including our own. Two constraint pitfalls we found
and fixed during this work had inflated earlier measurements.

For the field, we hope two ideas outlast the specific chains. First,
kernel choice (shared perimeter, density gradients, administrative
boundaries) can be a first-class modeling decision \emph{inside}
the proposal rather than a post-hoc filter, provided either the induced
distribution is corrected (Chain E) or the method is labeled plainly
as an optimizer (Chain C). Second, enumerated fibers are cheap,
severe test beds. A chain's distance from its own computable target
on 4\,006 states is a more informative diagnostic than any p-value on
a single trace. The most promising next step is the reversible
hybrid. It mixes Chain E's exact local kernel with a global tree move
Metropolized toward the same target via its computable
spanning-tree-count ratio, yielding a mobile chain whose stationary
distribution is known, the property that court-facing ensemble
analysis ultimately requires.

%════════════════════════════════════════════════════════════════════════════════
\appendix
\section{Algorithm details and hyperparameters}\label{app:hyperparams}
%════════════════════════════════════════════════════════════════════════════════

\subsection{Common setup}\label{app:common}

\begin{itemize}
  \item Dual graph $G = (V, E)$ with vertex populations $p_v$ and
    edge attributes (\texttt{shared\_perim} for Iowa).
  \item $k$ districts, $\varepsilon$-population balance per
    Equation~\eqref{eq:pop-balance}, enforced as a two-sided bound around the
    global ideal on every move (Equation~\eqref{eq:balance-window})
    and audited per step. The bare GerryChain \texttt{recom} proposal
    is wrapped with an explicit whole-plan balance check
    (\texttt{recom\_valid}), since its internal cut search validates
    only one side of each split.
  \item Initial partition from
    \texttt{gerrychain.tree.recursive\_tree\_part}, redrawn until all
    $k$ districts satisfy the global tolerance
    (\texttt{balanced\_tree\_part}).
  \item Step count: $5 \times 5$: $N = 1\,000$, burn-in $B = 100$;
    Iowa: $N = 500$, $B = 100$; long-run validation
    (\S\ref{sec:exp-longrun}): $2 \times 10^5$ (A) and
    $10^5$ (E) steps.
  \item Random seeds: base seed 42 with a distinct offset per chain
    (independent streams for Python \texttt{random} and
    \texttt{numpy}).
\end{itemize}

\subsection{Chain-specific parameters}\label{app:chains}

\begin{table}[H]
\centering
\small
\begin{tabular}{lll}
\toprule
Chain & Parameter & Value \\
\midrule
A & $\lambda$ (target temperature) & 0.10 \\
B & $p_R$ (ReCom mix probability) & 0.5 \\
B & $\alpha$ (Laplacian weight) & 10.0 \\
B & $\beta$ (fidelity weight) & 1.0 \\
B & Threshold scheme & Cheeger sweep, balance fallback \\
B & Laplacian normalization & False (unnormalized $L$) \\
C & $n_{\mathrm{iqp}}$ (IQP steps per outer iteration) & 10 \\
C & $\alpha$ & 10.0 \\
C & $\beta$ & 1.0 \\
C & MIQP time limit & 5.0 s \\
D & $p_R$ & 0.5 \\
D & Weight randomization & $\mathcal{U}[1, 2]$ per edge \\
D & Repair & False; balanced sweep in place of sign split \\
E & $\alpha$ & 1.0 \\
E & $\beta$ & 10.0 ($\tau = \alpha/\beta = 0.1$; Corollary~\ref{cor:small-tau} gives the resulting move size and \S\ref{sec:mobility} places the mobility optimum near $\tau = 0.3$) \\
E & Rounding & single-shot Bernoulli($x^\star$) \\
E & $\lambda$ (MH target) & 0.10 \\
\bottomrule
\end{tabular}
\caption{Locked-in hyperparameters for the six
  chains.}\label{tab:hyperparams}
\end{table}

\subsection{Diversity controls}\label{app:diversity}

Weight jitter multiplies each edge weight by a draw from
$\mathcal{U}[1, 2]$ before solving (always on for D, as in Davies et
al.'s Algorithm~2; never used in Chain E, whose detailed-balance
argument needs a deterministic QP map). Tabu forbids the last $K = 3$
pairs chosen by the local kernel, cleared whenever ReCom moves.
Their effects are reported in Appendix~\ref{sec:exp-ablations}.

\subsection{Compute environment}\label{app:compute}

\begin{itemize}
  \item macOS (Darwin 23.6.0), Python 3.12.
  \item Gurobi 13.0.2 (academic license) for QP and MIQP, at its
    default thread count.
  \item NetworkX 3.x for graph operations.
  \item One Python process per run; no GPU. All experiments were run
    on a 2024 MacBook.
\end{itemize}

%════════════════════════════════════════════════════════════════════════════════
\section{Partition fiber structure: communities without bottlenecks}\label{sec:fiber-corridors}

This subsection documents a structural diagnostic of the partition fiber
$\mathcal{P}_k(G; \varepsilon)$ (the set of all valid plans), viewed as a graph under its minimal
valid move relation. Two findings emerge, one positive and one
negative, and we report both. Positively, all three enumerated fibers
carry genuine mesoscale community structure, with modularity far above
degree-preserving null models, and chains differ sharply in how often
they cross community boundaries. Negatively, a plausible folk
intuition (that the communities are ``islands'' joined by narrow
corridors or bottleneck edges) is refuted by method-independent
statistics. Community interfaces are wide, nearly all community pairs
are directly connected, and the literal bridges of the exact-balance
fiber turn out to be pendant vertices. The diagnostic still motivates
the local-refinement design of \S\ref{sec:model}, and we believe both
findings are of independent interest for analyzing redistricting
MCMC.

\paragraph{Three test beds.}\label{sec:fiber-testbeds}

We construct three fully enumerated partition fibers spanning grid
size, district count, and balance tolerance.

\begin{table}[H]
\centering
\small
\begin{tabular}{lrrrrrl}
\toprule
Fiber & Nodes & Edges & $k$ & $\varepsilon$ & Sizes & Adjacency \\
\midrule
$5 \times 5$, $\varepsilon = 0$    & 4\,006   & 13\,416  & 5 & 0    & \{5\}       & boundary swap \\
$5 \times 5$, $\varepsilon = 0.20$ & 193\,128 & 969\,760 & 5 & 0.20 & \{4, 5, 6\} & single flip \\
$4 \times 4$, $\varepsilon = 0.25$ & 1\,953   & 7\,024   & 4 & 0.25 & \{3, 4, 5\} & single flip \\
\bottomrule
\end{tabular}
\caption{The three test-bed fibers. Adjacency is the \emph{minimal
  valid move} of each fiber. At $\varepsilon = 0$ no single-vertex
  reassignment preserves exact balance (it produces districts of sizes
  4 and 6), so adjacency is the boundary \emph{swap}, which exchanges one
  boundary vertex of district $a$ with one of district $b$, both
  districts remaining connected. At $\varepsilon > 0$ adjacency is the
  single \emph{flip}, which reassigns one boundary vertex, both affected
  districts remaining connected and within size bounds. Cross-fiber
  comparisons carry one caveat. The $\varepsilon = 0$
  fiber lives under a different (two-vertex) move relation than the
  $\varepsilon > 0$ fibers.}\label{tab:fibers}
\end{table}

For each fiber, all valid partitions are enumerated and the single-flip
adjacency graph is built explicitly (enumeration runtimes on a 2024
MacBook range from $\sim 5$ s for the $4 \times 4$,
$\varepsilon = 0.25$ fiber to $\sim 50$ min for the $5 \times 5$,
$\varepsilon = 0.20$ fiber; scripts accompany the replication
materials).

\paragraph{Community structure, validated against null models.}\label{sec:fiber-islands}

Each fiber graph carries substantial mesoscale community structure
under modularity maximization. We report the method per fiber (greedy
modularity, Clauset--Newman--Moore, for the $\varepsilon = 0$ and
$4 \times 4$ fibers; Louvain, resolution $1.0$, seed 42, for the large
$\varepsilon = 0.20$ fiber), validate against degree-preserving
rewiring nulls, and treat the specific community counts as
method-dependent descriptive statistics.

\begin{table}[H]
\centering
\small
\begin{tabular}{lrrrlrrr}
\toprule
& \multicolumn{4}{c}{Community structure}
& \multicolumn{3}{c}{Anti-bottleneck statistics} \\
\cmidrule(lr){2-5}\cmidrule(lr){6-8}
Fiber & Comm. & $Q$ & Null $Q$ & Method
      & Pairs conn. & Inter-comm.\ edges & Width (med/max) \\
\midrule
$5 \times 5$, $\varepsilon = 0$    & 10 & 0.65 & 0.36 & CNM
  & 15/15 & 18.3\% & 142 / 383 \\
$5 \times 5$, $\varepsilon = 0.20$ & 26 & 0.69 & n/a  & Louvain
  & 320/325 & 26.8\% & 667 / 3\,006 \\
$4 \times 4$, $\varepsilon = 0.25$ & 7  & 0.58 & 0.34 & CNM
  & 21/21 & 23.3\% & 71 / 186 \\
\bottomrule
\end{tabular}
\caption{Left: community decompositions with degree-preserving null
  baselines (mean $Q$ of CNM on double-edge-swap rewirings of the same
  graph; the $193$k-node fiber's null is omitted for cost). Real $Q$
  exceeds null $Q$ by $0.24$--$0.29$, so the community structure is
  genuine, not a degree-sequence artifact. Counts vary with the
  detection method and should be read qualitatively. Right:
  the statistics that refute the corridor picture. Nearly every pair
  of communities shares a direct interface (``pairs conn.''\ counts
  the large-community pairs at $\varepsilon = 0$), interfaces are
  tens to thousands of edges wide, and a fifth to a quarter of all
  fiber edges run between communities. The exact-balance fiber's 16
  literal single-edge bridges each split off a \emph{single vertex}
  (pendant leaves), not an island, and the six large
  $\varepsilon = 0$ communities have conductance $0.15$--$0.22$, an
  order of magnitude above bottleneck territory.}
\label{tab:communities}
\end{table}

\paragraph{No narrow corridors: a corrected intuition.}\label{sec:fiber-corridors-invariant}

It is tempting to picture these communities as islands joined by a few
narrow corridors, with bottleneck edges controlling passage. The
method-independent statistics on the right of
Table~\ref{tab:communities} refute that picture on all three
fibers.

Concretely, on the $\varepsilon = 0.20$ fiber, interfaces thinner than
50 edges do exist (there are 16 of them), but they are simply the tail
of a width distribution whose median is 667 and whose top interfaces
carry $3\,000$ edges; a chain moving through this graph crosses
between communities through the wide interfaces, and the thin tail is
irrelevant to mixing. The correct summary of fiber macro-structure is
therefore \emph{communities without bottlenecks}. Valid-plan space is
regionally organized but well connected, and slow traversal by local
chains (below) is a property of the \emph{dynamics}, not of corridor
topology. This negative result matters for chain design. It says
local chains stall not because the state space is nearly disconnected
but because their moves are small and, in Chain E's case, because
distant proposals are rarely accepted; global recombination moves, by
contrast, relocate across communities in a single step.

\paragraph{Chain step-trace diagnostic.}\label{sec:fiber-trace}

The community decomposition supports a direct diagnostic for any
redistricting chain. We label each visited partition by its
greedy-modularity community on the $\varepsilon = 0$ fiber and record
the \emph{jump rate}, the fraction of steps that move to a
different community.

\begin{table}[H]
\centering
\small
\begin{tabular}{lrrr}
\toprule
Chain & Communities visited & Jump rate & Jumps in 1\,000 steps \\
\midrule
V: Vanilla ReCom        & 8/10 & 15.3\% & 153 \\
A: Tempered ReCom       & 7/10 & 12.7\% & 127 \\
B: QP-diffusion (50/50) & 6/10 & 4.6\%  & 46  \\
C: ReCom + IQP local    & 6/10 & 10.5\% & 105 \\
D: ReCom + spectral mix & 7/10 & 9.3\%  & 93  \\
E: QP-MH exact          & 2/10 & 0.3\%  & 3   \\
\bottomrule
\end{tabular}
\caption{Inter-community jump rates on the $5 \times 5$,
  $\varepsilon = 0$ fiber (1\,000 steps per chain, seed 42, per-chain
  RNG streams). A community-blind baseline crosses at roughly the
  inter-community edge share (18.3\%,
  Table~\ref{tab:communities}); ReCom variants (V, A) sit near it
  at $\sim 13$--15\% per step, the local kernels (B, C, D) are
  stickier at $\sim 5$--11\%, and the strictly local QP-MH chain (E)
  almost never crosses, a property of its small accepted moves rather
  than of any bottleneck in the fiber.}\label{tab:island-trace}
\end{table}

\textbf{Reading of the diagnostic.} ReCom relocates across
communities in single steps; optimization-guided local kernels work
within them; and the exact-but-strictly-local Chain E shows the
failure mode of small accepted moves. (Per-chain step-traces, which
show the same contrast as dwell patterns over time, are released as
supplementary figures.) The regimes are complementary, and they
suggest an explicit design principle. A good redistricting sampler
should combine global moves that relocate across regions of plan
space with local moves that refine within them; ReCom does the
former, anchored local kernels the latter. This is the design of
Chains B and C, whose every step is a random mixture of a ReCom
(across-community) move and a diffusion (within-community) move. We
note the diagnostic's limits (one run per chain on one fiber, and
community labels from a specific detection method) and read it as
motivating evidence, not proof.

%════════════════════════════════════════════════════════════════════════════════
\section{Ablations and sensitivity}\label{sec:exp-ablations}

All tables in this appendix are from the corrected balance
constraint and Gurobi 13, on the same Chain C RNG stream as the
headline table.

\paragraph{Diversity controls (Jitter and Tabu).}\label{sec:exp-diversity}

\begin{table}[H]
\centering
\small
\begin{tabular}{lrrrr}
\toprule
Configuration & Mean cuts & Std & Distinct-plan rate & Wall-clock \\
\midrule
Baseline ($n_{\mathrm{iqp}} = 5$)               & 31.64 & 2.10 & 32.8\% & 31\,s \\
+ Jitter ($\mathcal{U}[1,2]$ on edge weights)   & 32.40 & 1.93 & \textbf{58.6\%} & 35\,s \\
+ Tabu (size 3)                                  & \textbf{30.16} & 1.60 & 14.6\% & 30\,s \\
\bottomrule
\end{tabular}
\caption{Diversity controls on Iowa ($n_{\mathrm{iqp}} = 5$,
  $\alpha = 10$, $\beta = 1$). ``Distinct-plan rate'' is the fraction
  of the 500 steps that produced a not-yet-seen plan. The two
  controls pull in opposite directions, as intended.}\label{tab:diversity-iowa}
\end{table}

The two controls are orthogonal knobs, and they pull in opposite
directions. Jitter multiplies each edge weight by a
$\mathcal{U}[1,2]$ factor before the solve, which breaks ties between
equally good cuts and sends the chain to different optima. It raises
the distinct-plan rate from $33\%$ to $59\%$ at a cost of under one
cut of compactness. Tabu forbids the last three merged pairs from
being reused, which forces each step to refine a boundary the chain
has not just touched. That deepens the optimization decisively (mean
$31.64 \to 30.16$, the best single configuration in this paper, with
std $2.10 \to 1.60$) while halving the distinct-plan rate, because
systematic coverage of the boundaries drives the chain into the same
deep local optima. In short, Jitter varies which cut is taken on a
given pair and buys exploration, while Tabu varies which pair is
picked and buys depth. A practitioner wanting a diverse ensemble
turns on Jitter; one wanting the sharpest compactness turns on Tabu.

\paragraph{Sensitivity to $n_{\mathrm{iqp}}$, $\alpha$, and $\beta$.}\label{sec:exp-sensitivity}

\begin{table}[H]
\centering
\small
\begin{tabular}{lrrr}
\toprule
Configuration & Mean cuts & $\Delta$ vs A & Wall-clock \\
\midrule
$n = 3, \alpha = 10, \beta = 1$  & 33.03 & $-16.9\%$ & 18\,s \\
$n = 5, \alpha = 10, \beta = 1$  & 31.64 & $-20.4\%$ & 31\,s \\
$n = 10, \alpha = 10, \beta = 1$ & 31.78 & $-20.1\%$ & 118\,s \\
$n = 3, \alpha = 50, \beta = 1$  & 32.70 & $-17.7\%$ & 24\,s \\
$n = 5, \alpha = 10, \beta = 0$  & 32.27 & $-18.8\%$ & 36\,s \\
$n = 10, \alpha = 50, \beta = 0$ & 31.19 & $-21.6\%$ & 196\,s \\
\bottomrule
\end{tabular}
\caption{Chain C parameter sensitivity on Iowa ($k = 4$,
  $\varepsilon = 0.05$, seed 42). $n_{\mathrm{iqp}}$ is the number of
  inner IQP passes per step. Wall-clock here wraps the whole run and
  was measured inside a 38-configuration sweep, whereas
  Table~\ref{tab:iowa-headline} sums per-step solve times on the same
  deterministic run, which is why the $n = 10$ row reads 118\,s here
  and 57\,s there. Compare rows within one table only.}\label{tab:sensitivity}
\end{table}

Table~\ref{tab:sensitivity} reports the uniform-kernel slice; the
full sweep over $n_{\mathrm{iqp}} \in \{3,5,10\}$,
$\alpha \in \{10,50\}$, $\beta \in \{0,1\}$ and all three kernels
(36 configurations) is released with the code as supplementary
material. Three readings. First, $n_{\mathrm{iqp}}$ has diminishing
returns.
Moving from $3$ to $5$ passes buys about three and a half percentage
points of cut reduction ($-17\%$ to $-20\%$), but moving from $5$ to
$10$ buys nothing more while multiplying the runtime by two to four
depending on the timing convention (see the caption), because after a
few passes the easy boundary improvements are exhausted and further
passes revisit already-refined pairs. Five passes is the practical
sweet spot. Second, $\alpha$ (the weight on the Laplacian smoothing
term) matters little once it is large enough to dominate. Moving from
$10$ to $50$ changes the result by half a cut or less, so we fix
$\alpha = 10$. Third, the
fidelity weight $\beta$ helps rather than hurts. At $n = 5$, keeping
$\beta = 1$ gives $-20.4\%$ against $-18.8\%$ with $\beta = 0$. The
anchor keeps each refinement near the current plan, which on
heterogeneous county populations lands better local optima than an
unanchored solve, and it is what the exactly reversible Chain E needs
to stay near-binary. We therefore keep $\beta = 1$ throughout.

%════════════════════════════════════════════════════════════════════════════════
\section{Reproducibility checklist}\label{app:reproducibility}
%════════════════════════════════════════════════════════════════════════════════

\begin{itemize}
  \item Code: \texttt{diffusion\_model/recom/} in the project
    repository.
  \item Modules: \texttt{problems.py} (graph setup, balanced initial
    plans), \texttt{chains.py} (chain runners, balance-validated
    ReCom wrapper, per-step audit), \texttt{qp\_model.py}
    (QP/IQP/spectral/QP-MH step factories, global-ideal balance
    windows), \texttt{diagnostics.py} (KL/TV diagnostics).
  \item Orchestrators: \texttt{iowa\_target\_matched.py} (mode, seed,
    and kernel as CLI arguments), \texttt{long\_run\_5x5.py}
    (\S\ref{sec:exp-longrun}), \texttt{ablation\_sweep.py}
    (\S\ref{sec:exp-ablations} and the supplementary sweep),
    \texttt{ensemble\_compare.py}, \texttt{sa\_baselines.py}, and
    \texttt{vanilla\_iqp\_cell.py}
    (\S\ref{sec:exp-optimizer-baselines}),
    \texttt{kernel\_traj\_metrics.py} and
    \texttt{multikernel\_figure.py} (\S\ref{sec:exp-multikernel}),
    \texttt{least\_change.py} (\S\ref{sec:exp-leastchange}),
    \texttt{matrix\_tree\_targets.py} (the Matrix--Tree
    quasi-target expectations of \S\ref{sec:exp-5x5}),
    \texttt{db\_check\_2x3.py} (the machine-precision kernel check of
    Remark~\ref{rem:qpmh-scope}),
    \texttt{ten\_seed\_stats.py} (between-seed statistics),
    \texttt{calibration\_distance.py} and \texttt{fig\_calibration.py}
    (the distributional calibration of \S\ref{sec:exp-longrun}, which
    reads only stored artifacts),
    \texttt{spectral\_move\_size.py} and \texttt{fig\_spectral.py}
    (the verification of \S\ref{sec:diffusion-theory}),
    \texttt{acceptance\_sweep.py} (Table~\ref{tab:mobility}), and
    \texttt{long\_run\_2x3.py} with \texttt{gof\_2x3.py} (the
    batch-means validation of \S\ref{sec:exp-2x3}; the long run is
    reproduced with arguments \texttt{1000000 42}, and the
    $5 \times 5$ long run of \S\ref{sec:exp-longrun} with
    \texttt{200000 100000 42}).
  \item Timing conventions differ between tables and are stated in
    each caption; Table~\ref{tab:iowa-headline} sums per-step solve
    times, Table~\ref{tab:sensitivity} wraps whole runs.
  \item Outputs:
    \texttt{iowa\_hybrid\_test\_out/target\_matched\_*\_seed*.npz},
    each stamped with the full configuration (kernel, $\alpha$,
    $\beta$, $n_{\mathrm{iqp}}$, balance mode, seed) and containing
    per-step cut traces, timings, and balance-audit arrays, so every
    table can be regenerated from its own artifact.
  \item Fibers: \texttt{fiber\_5x5/} contains the $5 \times 5$ and
    $4 \times 4$ enumeration scripts and precomputed
    adjacency files.
  \item Random seeds: base 42, per-chain offsets. Gurobi default seed
    (deterministic given the same input).
\end{itemize}

\bibliographystyle{plain}
\bibliography{ref}

\end{document}
